% !TEX root = ../PhD Thesis.tex
\chapter{The SG Transcriptome}
\label{chap:ChapterI}

\glspl{sg} continue to garner significant attention in various physiological and pathological contexts, owing to their assumed role in selective \gls{mrna} and protein sequestration, hypothesised to contribute to the downregulation of specific pathways in numerous diseases.
\\
\\
Two main methodologies have emerged for elucidating their \gls{rna} content, followed by standard \gls{rnaseq}: \gls{dc}, which in Khong \gls{etal} 2017 provided the first isolated \gls{sg} transcriptome\cite{khong_stress_2017}, and \gls{pl} via \gls{apexseq}\cite{padron_proximity_2019}, both relying on core \gls{sg} proteins (usually \gls{g3bp1}). 
\\
\\
DC relies on sequential centrifugation steps to fractionate cellular components based on their size and density\cite{khong_stress_2017}. Following cell lysis, low-speed spins remove nuclei and debris, while subsequent higher-speed centrifugation enriches for heavier cytoplasmic fractions containing large macromolecular assemblies, including SGs and other \gls{rnp} aggregates\cite{khong_stress_2017}. This approach enables the physical isolation of SG-enriched fractions, which can then be further purified, typically via immunoprecipitation of SG core proteins\cite{khong_stress_2017}.
\\
\\
In contrast, PL is based on the enzymatic tagging of biomolecules in the immediate vicinity of a protein of interest\cite{padron_proximity_2019}. In this approach, a bait protein, in this case a core SG protein, is fused to an engineered enzyme capable of covalently labelling nearby molecules (e.g., RNAs or proteins) within a defined spatial radius (around 20 \gls{nm})\cite{padron_proximity_2019}. Upon activation, the enzyme deposits a chemical tag (such as biotin) onto proximal targets, which can subsequently be purified using affinity-based methods\cite{padron_proximity_2019}. This strategy enables the identification of molecules that localise near the bait protein in living cells, without requiring physical isolation of the condensate itself\cite{padron_proximity_2019}.
\\
\\
Regardless, even when the same \gls{sg} core protein and stress induction are used, notable disparities between transcriptomes isolated with the different techniques are observed\cite{glauninger_stressful_2022}. 
\\
\\
While \gls{dc} highlights \gls{rna} length as a significant predictor of \gls{sg} enrichment\cite{khong_stress_2017}, this effect is absent with some \gls{pl} studies\cite{padron_proximity_2019,somasekharan_g3bp1-linked_2020}. The enrichment in longer \glspl{rna} has been hypothesised to play a functional role, as \glspl{sg} are formed by aggregation, longer transcripts equate to more potential binding sites for \glspl{rbp}, their main protein constituents \cite{khong_stress_2017}. Moreover, it is plausible that \gls{pl} also captures \gls{g3bp1} that does not localise to \glspl{sg}, and that therefore part of the obtained enrichment is not \gls{sg}-specific\cite{somasekharan_g3bp1-linked_2020}. 
\\
\\
In addition to transcript length, several other features have been proposed to influence \gls{sg} enrichment. These include translation efficiency and the number of binding sites for \gls{sg}-associated \glspl{rbp}\cite{khong_stress_2017}. However, since both features are themselves correlated with transcript length, such associations were expected\cite{glauninger_stressful_2022}. \gls{rna} modifications such as \gls{m6a}\cite{di_timoteo_m6a_2024} and \gls{ac4c}\cite{kudrin_n4-acetylcytidine_2024} have also been linked to \gls{sg} localisation. Yet, because the number of these modification sites also scales with transcript length, it remains controversial whether their enrichment persists after proper length normalization\cite{khong_limited_2022}.
In addition to transcript length, several other features have been proposed to influence \gls{sg} enrichment. These include translation efficiency and the number of binding sites for \gls{sg}-associated \glspl{rbp}\cite{khong_stress_2017}. However, since both features are themselves correlated with transcript length, such associations were expected\cite{glauninger_stressful_2022}. \gls{rna} modifications such as \gls{m6a}\cite{di_timoteo_m6a_2024} and \gls{ac4c}\cite{kudrin_n4-acetylcytidine_2024} have also been linked to \gls{sg} localisation. Yet, because the number of these modification sites also scales with transcript length, it remains controversial whether their enrichment persists after proper length normalisation\cite{khong_limited_2022}.
\\
\\
To address these discrepancies in \gls{sg} characterisation and to identify genuine transcript features associated with \gls{sg} enrichment, we performed a meta-analysis of publicly available human \gls{sg} transcriptomic datasets. These datasets span multiple cell lines, purification strategies, and stress conditions, and include studies employing \gls{dc}, \gls{pl}, as well as additional, less commonly used approaches such as \gls{pic}\cite{honda_high-depth_2021}, \gls{artrseq}\cite{xiao_profiling_2024}, \gls{eclipseq}\cite{sato_stress_2023}, \gls{fanci}\cite{zhou_rna_2024}, and \gls{photon}\cite{rajachandran_subcellular_2025}.

\section{Materials and Methods}
 
\subsection{Dataset Procurement and List}

For dataset procurement, we searched \gls{geo} Datasets using the term ``stress granule", filtering for expression profiling by high-throughput sequencing and restricting to human samples, excluding microarray studies. Using this approach, no single-cell datasets were retrieved. All publicly available studies up to September 1, 2024, were included. From the studies initially considered during the screening, we excluded specifically the dataset associated with Yi \gls{etal} (GSE245209 and GSE245210)\cite{yi_mapping_2024}, given that its \gls{coclip} method required peak-based analyses across conditions, more analogous to \gls{clip} than standard differential expression. Part of the study from Sato \gls{etal} 2023\cite{sato_stress_2023} also featured samples more akin to \gls{clip}, but, as these had considerably higher read depth, we found they could be analysed as regular \gls{rnaseq} samples, and were thus kept.
\\
\linebreak
The following datasets identified by their \gls{gseid} (more detailed in Table 3.1), and additional data/supplemental information, divided by their usage, were used throughout this work:
\begin{itemize}
	\item For the \gls{sg} transcriptomes: GSE99304\cite{khong_stress_2017}; GSE138988\cite{matheny_transcriptome-wide_2019}; GSE121575\cite{padron_proximity_2019}; GSE138058\cite{somasekharan_g3bp1-linked_2020}; GSE119977\cite{matheny_rna_2021, matheny_corrigendum_2023}; GSE143413\cite{honda_high-depth_2021}; GSE171792\cite{shen_sexually_2022}; GSE188397\cite{sato_stress_2023}; GSE171655\cite{iadevaia_novel_2022}; GSE197001\cite{curdy_proteome_2023}; GSE223295\cite{ren_profiling_2023}; GSE226161\cite{xiao_profiling_2024}; GSE218180\cite{zhou_rna_2024}; GSE212380\cite{kudrin_n4-acetylcytidine_2024}; GSE242766\cite{di_timoteo_m6a_2024}; PRJNA115262153\cite{rajachandran_subcellular_2025}.
	\item For the \gls{pb} transcriptomes: GSE13898834\cite{matheny_transcriptome-wide_2019}; GSE224858/GSE224752\cite{kodali_rna_2024}.
	\item For yeast centrifugation modulation: GSE265963\cite{glauninger_transcriptome-wide_2025}.
	\item Canonical transcript lengths were obtained from ENSEMBL\cite{cunningham_ensembl_2022}.
	\item \gls{ac4c} \gls{rna} modification information was obtained from Arango \gls{etal} 2022\cite{arango_direct_2022}'s Supplementary Table 2, from the ``ac4C\_gene\_summary" sheet.
	\item \gls{m6a} \gls{rna} modification information was obtained from Khong \gls{etal} 2022's Supplementary Data 1.
\end{itemize}

\subsection{Correlations}
 
Unless otherwise indicated, correlations were calculated by comparing gene-level moderated t-statistics obtained from differential expression analyses, with Spearman’s correlation coefficient $\rho$ and the associated test’s p-value. For correlations involving gene length, the canonical transcript length (in number of \glspl{nt}), obtained from ENSEMBL\cite{cunningham_ensembl_2022} was used, with the longest transcript chosen in cases of ambiguity.

\subsection{Sampling Bias Simulations (Mimi-Seq)}

We first considered that a technical sampling bias could explain the enrichment of longer \gls{rna}s in \gls{dc}-isolated \gls{sg}s. In this context, technical sampling bias arises from the fact that RNA-seq measures transcripts through randomly generated fragments\cite{chu_rna_2012}. For a given expression level, longer transcripts produce a greater number of fragments simply due to their increased length, leading to a higher probability of detection and quantification\cite{chu_rna_2012}. This effect can be further compounded by size-dependent biases during library preparation, fragmentation, and sequencing\cite{chu_rna_2012}. As a result, longer transcripts can appear artificially enriched, independently of true biological localisation.
\\
\\
To assess this, we simulated RNA-seq fragmentation using expression levels and fragment length distributions derived from GSE99304\cite{khong_stress_2017} \gls{wc} and \gls{sg} samples. PE reads were first aligned to the reference transcriptome (\gls{v} section 3.2.1), and the resulting \gls{bam} files were used to extract two key metrics for each sample:
\begin{enumerate}
	\item the observed \gls{tlen} field, which reports the observed fragment size;
	\item the median fragment position in the transcript, calculated for each read pair as:
	\begin{equation}
		\frac{\text{start}(\text{read}_1) + \text{start}(\text{read}_2) + \text{read length}}{2}
	\end{equation}
	\eqcaption{Median fragment position}
\end{enumerate}
\noindent
For this dataset, read length was 152 \gls{nt}. Fragment length distributions were represented as empirical probability tables describing the frequency of each observed fragment size. Transcript abundances were first normalised to \gls{tpm}, providing length-normalised estimates of transcript representation independent of transcript size and library depth. TPM values were then scaled according to the target sequencing depth of either \gls{wc} (1,000,000 reads) or \gls{sg} (100,000 reads) libraries, reflecting reports that \gls{sg}s contain approximately 10\% of total cellular \gls{rna}\cite{khong_stress_2017}. These scaled abundances were used to estimate the number of fragmentation events assigned to each transcript.
\\
\\
For each fragmentation event, fragment lengths were randomly sampled from the corresponding empirical size distribution, and transcripts were iteratively fragmented until no additional fragment could be accommodated within the remaining sequence. This approach preserved both transcript identity and fragment length characteristics in the simulated reads.
\\
\\
Following simulation, transcript abundances were converted back to discrete read counts to generate sequencing-like count matrices, allowing downstream differential expression analyses, using the same pipeline as described in section 3.2, while preserving the intended total library depth of empirical \gls{sg} datasets. Three simulated WC and three simulated SG samples were generated to assess whether subsampling to approximately 10\% of total reads, mimicking the SG fraction, was sufficient to induce length-dependent transcript representation biases.


\subsection{Centrifugation Simulation (Spin-Sim)}

We then considered that a physical centrifugation bias could explain the enrichment of longer \gls{rna}s in \gls{dc}-isolated \gls{sg}s, as longer, and therefore heavier, transcript species are expected to more readily precipitate\cite{glauninger_transcriptome-wide_2025}.
\\
\\
In this section, the term \gls{rna}, and its associated \gls{mw}, is used operationally to refer to the average enrichment or sedimentation behaviour associated with a given transcript species. This reflects the collective behaviour of all molecules of that transcript across the population of aggregates, rather than individual \gls{rna} molecules, which do not precipitate independently but instead co-sediment as part of larger ribonucleoprotein condensates\cite{glauninger_transcriptome-wide_2025}.
\\
\\
Read counts were first converted to \gls{tpm}, providing length-normalised estimates of transcript representation independent of transcript size and library depth. Reference nucleotide sequences for each transcript were retrieved from ENSEMBL\cite{cunningham_ensembl_2022} to determine nucleotide composition and calculate the molecular weight (\gls{mw}) of each \gls{rna}. Molecular weights were computed as the sum of the individual nucleotide masses, based on standard values\cite{thermo_fisher_dna_nodate}, following IUPAC nomenclature\cite{noauthor_nomenclature_1985}:

\begin{equation}
	MW = (A \times 329.2) + (U \times 306.2) + (C \times 305.2) + (G \times 345.2) + 159 \;\text{g mol}^{-1}
\end{equation}
\eqcaption{RNA Molecular Weight}
\\
For \gls{rna} modification experiments, transcript \gls{mw} was adjusted by $+14\,\mathrm{g\,mol^{-1}}$ for each m6A site\cite{national_center_for_biotechnology_information_pubchem_2025} and $+42\,\mathrm{g\,mol^{-1}}$ for each ac4C site\cite{national_center_for_biotechnology_information_pubchem_2025-1}.
\\
\linebreak
\noindent
To simulate RNA sedimentation during centrifugation, we applied a simplified particle sedimentation model\cite{svedberg_ultracentrifuge_1940}: 
\begin{equation}
	\text{Distance} =
	\text{Sedimentation coefficient} \times
	\text{Acceleration} \times
	\text{Time}
\end{equation}
\eqcaption{RNA Sedimentation Distance}
\\
Since precise sedimentation coefficients depend on transcript-specific features (e.g., shape, folding), we used a heuristic approximation\cite{balch_effects_2015} assuming spherical \gls{rna} behaviour and sedimentation coefficients proportional to the cube root of \gls{mw}: 
\begin{equation}
	\text{Distance} = \sqrt[3]{\text{MW}} \times \text{Acceleration} \times \text{Time}
\end{equation}
\eqcaption{RNA Sedimentation Heuristic Distance}
\\
The centrifugation process was modelled after Khong \gls{etal} 2017\cite{khong_stress_2017}, assuming initiation from the 1 mL mark in standard 2 mL microcentrifuge tubes (resulting in 2 cm from the bottom, as empirically measured using standard safe-lock tubes). Each transcript’s sedimentation distance was scaled proportionally to estimate the fraction of each RNA that would sediment. These sedimented proportions were converted back to read counts, to generate sequencing-like count matrices as above-mentioned, using:
\begin{equation}
	\text{Counts} =
	\left(
	\frac{\mathrm{TPM} \times \text{Transcript Length}}{10^{6}}
	\right)
	\times \text{Original Read Depth}
\end{equation}
\eqcaption{TPMs to Counts Conversion}
\noindent
where Original Read Depth corresponds to the total number of reads in the associated \gls{sg} sample, ensuring consistency between original and simulated datasets.
\\
\linebreak
\noindent
Additional simulations used empirical data from Glauninger \gls{etal} 2025\cite{glauninger_transcriptome-wide_2025}, which measured \gls{rna} sedimentation in centrifuged yeast samples. From these dataset’s sodium azide stressed samples, which we considered was the most similar to the commonly used arsenite, as they both cause oxidative stress\cite{garg_elucidation_2020}, we derived a linear regression model relating \gls{rna} length to sedimentation coefficients. Specifically, an ordinary least squares model was implemented in R using the \texttt{lm()} function, allowing estimation of the linear relationship between transcript length and sedimentation behaviour. This model showed a slightly stronger association than an equivalent model using \gls{mw} as predictor. After testing different models, we used the cube root of RNA length to fit the final model:
\begin{equation}
	\text{Sedimentation Coefficient} =
	6.09 \times 10^{-11}\sqrt[3]{\text{RNA Length}} + 3.88 \times 10^{-10}
\end{equation}
\eqcaption{Empirical Derived RNA Sedimentation Coefficients}
\\
Mitochondrial \glspl{rna} exhibited disproportionately high sedimentation relative to size, so we derived a separate equation for them using the same method as above, but considering only mitochondrial encoded \glspl{rna}:
\begin{equation}
	\text{MT Sedimentation Coefficient} =
	3.23 \times 10^{-11} \sqrt[3]{\text{RNA Length}} + 1.17 \times 10^{-9}
\end{equation}
\eqcaption{Empirical Derived Mitochondrial RNA Sedimentation Coefficients}
\\
As before, sedimented \gls{rna} abundances were converted back to read counts using the same formula. 

\subsection{\textit{In Silico} Resources}

All analyses were performed on the workstation \textit{lobito} at the \gls{gimm} data centre, which is equipped with 19 \gls{tb} of \gls{ssd} storage, 768 \gls{gb} of RAM, 72 \gls{cpu} threads (\gls{intel} Xeon Gold 6254 @ 3.10 \gls{ghz}), and an integrated Matrox G200eW3 \gls{gpu}. \gls{bash} was used to run FASTQC\cite{andrews_fastqc_2010}, MultiQC\cite{ewels_multiqc_2016}, STAR\cite{dobin_star_2013}, featureCounts\cite{liao_featurecounts_2014}, and IGV visualisation. Downstream analyses were conducted in R (v4.1.2) using RStudio Server (v1.4.1717). Mimi-Seq and SpinSim simulations were conducted in Python 3 (v3.10.11). 

\section{Standard Bioinformatics Analytical Pipeline}

The following section summarises the core analyses consistently applied throughout this thesis, together with representative examples illustrating their interpretation and practical importance. This section is intended as a guide, typically highlighting both a ``good" and a ``less good" example for each type of analysis.

\subsection{Quality Control: Dataset Pre-Processing}

Before conducting any analyses, it is essential to evaluate the quality of the \gls{rnaseq} data, as technical issues associated with library preparation and sequencing can strongly influence the results. Tools such as FastQC\cite{andrews_fastqc_2010} and MultiQC\cite{ewels_multiqc_2016} allow inspection of the raw FASTQ files and provide several metrics related to read quality and sequence composition (Figure 3.1). 
\\
\\
For example, FastQC reports per-base quality scores across sequencing reads, enabling the identification of systematic declines in sequencing quality\cite{andrews_fastqc_2010}. In \gls{rnaseq} datasets, reduced base-calling accuracy is commonly observed toward the 3$'$ end of reads, accumulating over successive sequencing cycles for two principal reasons: incomplete removal of fluorescently labelled nucleotides between cycles introduces carry-over signal; and phasing effects, whereby individual molecules within a cluster progressively fall out of synchrony, cause the detected fluorescence to become an increasingly mixed signal. In \gls{pe} sequencing, the second read of a pair typically exhibits slightly lower quality than the first, as it is generated after an additional round of sequencing chemistry and therefore accumulates a greater degree of both effects (Figure 3.1A).
\\
\\
FastQC also identifies overrepresented sequences, which can correspond to sequencing adapters if adapter trimming was not performed, or to \gls{rrna} fragments if ribosomal depletion during library preparation was inefficient\cite{andrews_fastqc_2010} (Figure 3.1B). 
\\
\\
Alignment or read summarisation statistics provide an additional layer of quality assessment (Figure 3.1C). In general, a uniquely mapped read percentage above ~70\% is considered acceptable for \gls{rnaseq} experiments\cite{peter_fitzgerald_introduction_2024}. Lower mapping rates may indicate degraded \gls{rna}, contamination from another organism, or an excess of repetitive sequences such as \gls{rrna} resulting from incomplete ribosomal depletion. 
\begin{figure}[H]
	\centering
	\includegraphics[width=1\linewidth]{"images/Paper Images/TESE_Fig4.1"}
	\caption[Different types of information provided by FastQC]{\textbf{Different types of information provided by FASTQC.} (A) Per base sequence quality in paired-end reads, in first (Left) and in second (Right) pair members, from sample SRR5605168 (GSE99304\cite{khong_stress_2017}). Note the lower scores on the second member of the read pair. Image is a screenshot from FastQC's resulting \gls{html} files. (B) Overrepresented sequences in FASTQ files, prior to adapter trimming in (Left) SRR10191283\_1 (GSE138058\cite{somasekharan_g3bp1-linked_2020}) and in (Right) SRR8083956 (GSE121575\cite{padron_proximity_2019}). Sequences from Left represent a higher total percentage of reads, individually, and originate from \gls{rnaseq} adapters. Sequences from Right represent a lower total percentage of reads, individually, and originate from rRNAs, as ribosomal depletion was intentionally skipped in this study\cite{padron_proximity_2019}. Image is a screenshot from FastQC's resulting \gls{html} files. (C) featureCounts\cite{liao_featurecounts_2014} summary of assigned and unassigned reads, from the same files represented in B, after adapter trimming. In Left, a suitable number of reads were uniquely mapped. In Right, a very small percentage of reads was uniquely mapped, as ribosomal depletion was intentionally skipped in this study\cite{padron_proximity_2019}.}
	\label{fig:fastqc}
\end{figure}
\noindent
Throughout this thesis, raw sequencing reads (FASTQ format) were preprocessed using Trimmomatic (v0.39)\cite{bolger_trimmomatic_2014}, operating in either \gls{se} or \gls{pe} mode depending on the original study design. Trimming was performed to remove sequencing adapters, low-quality bases, and technical artefacts arising during library preparation and sequencing, which can otherwise compromise alignment accuracy. Appropriate adapter sequences were used for each dataset, and quality filtering parameters were applied to retain high-confidence read segments, using default parameters.
\\
\\
Trimmed reads were then aligned to the \gls{hg38} reference genome using \gls{star} (v2.7.9a)\cite{dobin_star_2013} with custom parameters designed to prioritise alignment specificity. In particular, only uniquely mapped reads were retained (--outFilterMultimapNmax 1), thereby excluding reads that map equally well to multiple genomic loci. Additional constraints on mismatch rates (--outFilterMismatchNoverLmax 0.04) and splice junction support (--alignSJoverhangMin 8) further reduced ambiguous or low-confidence alignments. Collectively, these parameters favour high-confidence read assignment at the cost of reduced sensitivity, as reads with multiple possible mappings or higher mismatch rates are discarded\footnote{Although less stringent parameters would increase overall read assignment in datasets such as GSE121575\cite{padron_proximity_2019}, where ribosomal depletion was not performed, the objective of this \gls{phd} project was to characterise the \gls{sg} transcriptome with high specificity, motivating the exclusion of ambiguously mapped reads.}. To ensure consistency across datasets, the following parameters were applied in all alignments:

\begin{lstlisting}
	--outFilterType BySJout
	--alignSJoverhangMin 8
	--alignSJDBoverhangMin 1
	--alignIntronMin 20
	--alignIntronMax 1000000
	--alignMatesGapMax 1000000
	--outFilterMultimapNmax 1
	--outFilterMismatchNmax 999
	--outFilterMismatchNoverLmax 0.04
\end{lstlisting}

\noindent
Aligned reads (\gls{bam} files) were subsequently annotated using the GENCODE v42 primary assembly\cite{frankish_gencode_2021} and quantified with featureCounts (v2.0.3)\cite{liao_featurecounts_2014}.

\subsection{Quality Control: Coverage and Complexity}

After these initial quality control steps and read counting, two additional parameters are important to assess at the count matrix level: sample coverage and sample complexity. 
\begin{itemize}
\item \textbf{Sample coverage} (Figure 3.2A) reflects the effective sequencing depth, defined as the total number of reads assigned to annotated genes in a given sample\cite{costa_uncovering_2010}. To visualise this, cumulative coverage plots were generated by ranking genes by decreasing expression level and computing the cumulative number of reads as progressively more genes are included. The X-axis represents the cumulative number of genes ordered from highest to lowest expression, while the Y-axis shows the cumulative fraction of total reads. Typically, the curve rises steeply for the most highly expressed genes before gradually flattening, reflecting the log-normal distribution of gene expression in which a small number of transcripts dominates the total read count. Notably, differences in the shape and rate of this flattening across samples are informative of differences in library complexity, with steeper curves indicating greater domination by a small subset of transcripts.

\item \textbf{Sample complexity} (Figure 3.2B) describes the distribution of reads across genes, providing a measure of library diversity\cite{costa_uncovering_2010}. It was estimated as the number of genes required to account for \textit{Y}\% of total reads, computed by ranking genes by decreasing expression and calculating their cumulative contribution until the full library was accounted for. Samples with low complexity reach high read proportions with relatively few genes, indicating that a small subset of highly expressed transcripts dominates the library; conversely, high-complexity samples require a greater number of genes to reach the same proportions, reflecting a more even read distribution. Complexity is strongly influenced by the total amount of mRNA present in the input sample. Libraries prepared from greater quantities of mRNA capture a broader representation of the transcriptome's diversity, whereas libraries from low-input samples are more susceptible to stochastic sampling of only the most abundant species. Additionally, technical factors such as incomplete rRNA depletion can further reduce complexity, as highly abundant ribosomal transcripts then dominate the sequenced reads at the expense of mRNA representation\cite{costa_uncovering_2010}.
\end{itemize}
These metrics are particularly useful for identifying outlier or potentially compromised samples. In this context, problematic samples are defined as those displaying substantial deviations in coverage or complexity relative to the cohort, indicative of technical issues such as low sequencing depth, poor mapping efficiency (reflecting a low fraction of reads aligning to the reference genome), or overrepresentation of a small subset of transcripts.
\\
\\
In cases where these deviations are extreme, (i.e., markedly reduced total assigned reads or very low complexity, where a small number of genes account for a large fraction of the library) the sample was considered technically compromised and excluded from downstream analyses. More moderate deviations, however, were not grounds for immediate exclusion, as differences in coverage and complexity can often be partially mitigated through normalisation procedures\cite{costa_uncovering_2010}.
\\
\\
To further assess sample quality, exploratory analyses such as \gls{pca} were performed on normalised log-transformed gene expression (\gls{cpm}). Samples that assemble separately from their biological group were examined in conjunction with their coverage and complexity profiles. When anomalous behaviour in \gls{pca} coincided with clear technical deficiencies (low coverage and/or reduced complexity), this was interpreted as evidence of a technical artefact, and the sample was excluded. In contrast, samples with atypical clustering but without evident technical issues were retained, as such variation may reflect underlying biological differences rather than data quality concerns.
\\
\\
All of these assessments were performed in an empirical manner, on a dataset by dataset basis. 

\begin{figure}[H]
	\centering
	\includegraphics[width=1\linewidth]{"images/Paper Images/TESE_Fig4.2"}
	\caption[Sample Coverage and Complexity]{\textbf{Sample coverage and complexity.} (A) Cumulative read coverage plot showing the distribution of sequencing reads across genes for each sample (i.e., how many reads are assigned to how many genes). The data correspond to samples from GSE99304\cite{khong_stress_2017}. In this example, differences in coverage appear to be sample-type specific, as \gls{sg} samples consistently exhibit higher coverage than their \gls{wc} counterparts. (B) Sample complexity plot showing the number of genes required to account for \textit{X}\% of the total reads. The data correspond to samples from GSE143413\cite{honda_high-depth_2021}. In this case, two samples display markedly lower complexity, clearly observed at around 75\% of the total reads, represented by a dashed line. As these represent technical control samples in the original study\cite{honda_high-depth_2021}, this pattern is expected. However, it is important to note that direct comparisons involving these samples will inevitably yield numerous apparent differences due solely to their lower complexity.}
	\label{fig:CovandComp}
\end{figure}

\subsection{Processing Steps: Filtering and Normalisation}

Following the initial quality control procedures, the next step in the analysis pipeline is gene \textbf{filtering}. Although this step is often not stringent enough \cite{padron_proximity_2019, shen_sexually_2022}, it is essential for improving the reliability of downstream analyses. Filtering removes genes with extremely low expression levels, retaining only those that are likely to provide meaningful biological information. Low-count genes contribute little statistical power in differential expression analyses and can inflate variance estimates, thereby reducing the sensitivity of statistical tests\cite{costa_uncovering_2010}.
\\
\linebreak
\noindent
To illustrate why filtering is important, consider the following simplified example:
\begin{verbatim}
Gene A: 1000 reads in condition X and 2000 reads in condition Y
Gene B: 2 reads in condition X and 10 reads in condition Y
\end{verbatim}
Gene B displays a larger fold change (five-fold) than Gene A (two-fold). However, the biological relevance of Gene B’s change is questionable. A difference of only eight reads may simply reflect stochastic sampling noise rather than a meaningful biological signal. In contrast, the increase of one thousand reads observed for Gene A represents a substantial change in transcriptional activity and is therefore far more likely to reflect a genuine biological effect. Filtering helps prevent low-count genes such as Gene B from disproportionately influencing the statistical analysis\footnote{This issue appears to occur in GSE121575\cite{padron_proximity_2019}. Several genes reported in the original paper as differentially expressed with logFC $>$ 10 correspond to changes from 0 to approximately 10 reads, suggesting that insufficient filtering of low-count genes may have inflated the apparent fold changes. A similar issue is observed in the proteomic section of the same study, where one of the proteins reported as enriched in \glspl{sg} is \gls{ddx3y}, a Y-chromosome–encoded protein. All experiments in that study were performed in \gls{hek293t} cells, a cell line derived from a biological female with a 46,XX karyotype, making the presence of a Y-linked protein highly unlikely and thus indicative of technical or analytical artefacts.}.
\\
\\
Filtering thresholds are typically determined by visualising the distribution of read counts across genes (Figure 3.3A). In many \gls{rnaseq} datasets, this distribution appears bimodal, with one peak corresponding to genes with extremely low counts and another representing genuinely expressed genes. The low-count peak often arises from background noise, sequencing artifacts, or sporadic transcriptional events. By selecting a threshold between these two modes, genes with negligible expression can be removed from the analysis, which otherwise could lead to misleading interpretation (Figure 3.3B). In this thesis, these were defined not as hard thresholds, but on a dataset-by-dataset basis.
\\
\\
After filtering, the next critical step is \textbf{normalisation}. Normalisation corrects for differences in coverage and, to some extent, complexity between samples. Without normalisation, comparisons of gene expression across samples may be misleading.
\\
\linebreak
\noindent
Consider the following simplified example:
\begin{verbatim}
Gene C: 1000 reads in Condition 1 and 1000 reads in Condition 2
Gene D: 500 reads in Condition 1 and 500 reads in Condition 2
\end{verbatim}
At first glance, the expression levels appear identical between conditions. However, suppose the total sequencing coverage differs substantially: Samples in condition 1 were sequenced with a read depth of 5 million total reads, whereas Condition 2 with 20 million total reads. In this scenario, the same raw read counts represent a much larger fraction of the transcriptome in Condition 1 than in Condition 2. Briefly, normalisation adjusts the counts by dividing them by the total number of read counts in that sample, enabling expression levels to be compared on a common scale. As aforementioned, appropriate filtering prior to normalisation is essential. Widely used \gls{rnaseq} normalisation methods, including the \gls{tmm} approach implemented in \texttt{edgeR}\cite{robinson_edger_2010} and used throughout this thesis, estimate scaling factors from a large set of genes that are robustly expressed, unlikely to dominate the library, and expected to remain stable across conditions, serving as a reliable centroid of transcriptomic behaviour from which sample-wide scaling can be confidently derived\cite{robinson_edger_2010,robinson_scaling_2010}. Lowly expressed genes are particularly susceptible to sampling ``noise" and exhibit unstable expression estimates, contributing disproportionately to technical variability relative to their biological signal\cite{robinson_edger_2010,robinson_scaling_2010}. If retained in large numbers, these genes can bias the estimation of library scaling factors and reduce the effectiveness of downstream normalisation\cite{robinson_edger_2010,robinson_scaling_2010}. Pre-filtering therefore removes genes with insufficient or unreliable signal, ensuring that normalisation is driven primarily by robustly detected transcripts with more reliable expression measurements\cite{robinson_edger_2010,robinson_scaling_2010}.
\\
\\
In this work, normalisation factors were estimated using the \gls{tmm} method\cite{robinson_scaling_2010} implemented in the R package ``\gls{edger}"\cite{robinson_edger_2010}. Unlike simple library-size scaling, \gls{tmm} additionally corrects for composition biases arising from the disproportionate contribution of highly expressed transcripts to total library counts\cite{robinson_scaling_2010}. Following filtering, read counts were normalised using the calculated \gls{tmm} scaling factors and subsequently transformed to log$_2$ \gls{cpm} using the ``\gls{voom}"\cite{law_voom_2014} function from the R package ``\gls{limma}"\cite{ritchie_limma_2015}, which models the mean-variance relationship of \gls{rnaseq} data for downstream linear modelling (Figure 3.3C).

\begin{figure}[H]
	\centering
	\includegraphics[width=1\linewidth]{"images/Paper Images/TESE_Fig4.3"}
	\caption[Filtering and normalisation]{\textbf{Filtering and Normalisation.} (A) Density plot showing the distribution of summed reads per gene (log10 scale) before (Left) and after (Right) filtering, in samples from GSE188397\cite{sato_stress_2023}. In this example, the filtering threshold was set at log10 = 3.25, corresponding to the empirical boundary between the lowly-expressed and other gene populations. (B) Scatter plot comparing gene expression in two replicates from the same sample type in GSE171792\cite{shen_sexually_2022}; genes as circles, axes given in log2 \gls{cpm}, Spearman’s correlation coefficient $\rho$ and associated test’s p-value on top left corner. After filtering (Right), the samples appear less similar than before filtering (Left). This occurs because many genes with zero or near-zero counts artificially inflate correlation estimates when retained. Even after filtering, the similarity between replicates remains somewhat lower than reported in the original study\cite{shen_sexually_2022}, possibly due to differences in normalisation methods. (C) Box plots showing the distribution of read counts across samples from PRJNA115262153\cite{rajachandran_subcellular_2025} before (left) and after (right) normalisation. Prior to normalisation, noticeable differences between samples are present, reflecting variation in sequencing depth and library composition. These discrepancies are largely resolved after the normalisation step. Colours correspond to sample types, as per the original study\cite{rajachandran_subcellular_2025}, given beneath the plots. Although this distinction is not relevant for the purpose of this example, it was kept to illustrate that read count distribution is often sample type related. In this case, samples correspond to \gls{photon}-isolated nuclei (Nucleus), and SGs in \gls{m6a} \gls{wt} (SG) and \gls{ko} (SG \gls{m6a} \gls{ko}), as well as their respective inputs (corresponding total cellular \gls{rna}\cite{rajachandran_subcellular_2025}).}
	\label{fig:FiltandNorm}
\end{figure}

\subsection{Exploratory Analysis: PCA}

A common first step in transcriptomic data exploration is \gls{pca}. \gls{pca} is a dimensionality reduction technique that summarises large datasets by identifying the main sources of variation across samples. Rather than examining thousands of genes individually, \gls{pca} reduces the data into a small number of \glspl{pc} that capture the largest patterns of variability within the dataset. In this case, \gls{pca} was performed on normalised, log2-transformed expression values to stabilise variance across genes, using mean-centering without variance scaling.
\\
\\
In well-defined biological scenarios, samples often position according to experimental conditions in \gls{pca} space. Such separation indicates that the biological phenotype under investigation is a major contributor to transcriptomic variation (Figure 3.4A).
\\
\\
Because \gls{pca} plots typically display only \gls{pc}1 \gls{vs} \gls{pc}2, additional principal components may still capture relevant sources of variation. To evaluate this possibility, it is useful to generate a scree plot (Figure 3.4B), which shows the proportion of variance explained by each principal component. This allows one to assess whether other components should also be explored visually, for example by plotting \gls{pc}2 \gls{vs} \gls{pc}3.
\\
\\
\Gls{pca} plots can reveal issues affecting specific samples, including technical variation associated with sample complexity or sequencing coverage, sample label misattribution (Figure 3.4C), or replicate samples displaying unexpectedly high transcriptional divergence (Figure 3.4D).
\\
\\
The absence of clear separation in \gls{pca} space, or along specific principal components, suggests that the experimental condition of interest does not account for a major source of variance in the dataset (Figure 3.4E). In other words, the variance captured by the leading principal components is dominated by other factors, such as technical variation or unrelated biological differences. Importantly, this does not imply that biologically meaningful effects are absent. Rather, it indicates that condition-associated transcriptional changes are likely modest in magnitude or restricted to a subset of genes, and therefore not reflected in the dominant axes of variation.
\\
\\
As an additional layer of interpretation, principal components can be examined through their gene loadings, which quantify the contribution of each gene to a given component. In other words, one can determine which genes are the primary drivers of the observed variance. In the example from GSE242766\cite{di_timoteo_m6a_2024} (Figure 3.4F), many of the genes contributing most strongly to \gls{pc}1, and to a lesser extent \gls{pc}2, correspond to transcripts at the extremes of gene length, either very short or very long. This is consistent with the observation that sample type, either \gls{sg} or \gls{wc}, appears to explain these components and with the proposed role of gene length in SG enrichment\cite{khong_stress_2017}.

\begin{figure}[H]
	\centering
	\includegraphics[width=1\linewidth]{"images/Paper Images/TESE_Fig4.4"}
	\caption[PCA as an Exploratory Tool]{}
	\label{fig:PCAExamples}
\end{figure}

\clearpage

\begin{figure}[t]
	\captionsetup{labelformat=adja-page}
	\ContinuedFloat
	\caption[]{\textbf{PCA as an exploratory tool.} (A) \gls{pca} plot of gene expression GSE99304\cite{khong_stress_2017} samples (\gls{u2os} cells). Apparent distances in \gls{pca} space can sometimes be misleading. Although a few samples appear separated, this separation occurs along \gls{pc}2, which explains only 2.94\% of the variance. In contrast, most of the variance in this dataset is explained by \gls{pc}1 (95.27\%), which corresponds to sample type (i.e., if the transcriptomes are from \gls{sg} or \gls{wc}). (B) Scree plot of the same PCA of panel A. In this case, the ``scree" occurs immediately after \gls{pc}1, indicating that the most meaningful dataset's variance is captured by this component. (C) PCA plot of gene expression in  GSE119977\cite{matheny_rna_2021} samples (\gls{u2os} cells). Samples are coloured in accordance to their G3BP1 state (i.e., if the protein is \gls{ko} (in red), conjugated with \gls{gfp} (in green), or \gls{wt} (in blue)\cite{matheny_rna_2021}). In the sample groups identified by crimsom arrows, the \gls{ko} and \gls{wt} samples within, differ only in this metadata variable. However, a pair of G3BP1 \gls{ko} samples cluster with their \gls{wt} counterparts in each group identified by crimsom arrows, suggesting possible sample mislabelling\protect\footnotemark. (D) PCA plot of gene expression in GSE171792\cite{shen_sexually_2022} samples (\gls{hela} cells). Samples are coloured in accordance to \gls{pl} target (GFP as control (green), DDX3X (red), or DDX3Y (blue)), and stress state (unstressed (lighter shades) or stressed with arsenite (darker shades)). In this case, the dark green samples do not cluster together, with one in particular deviating markedly from the remaining samples, to a degree exceeding the separation observed between the GFP Stressed condition and the other groups. This suggests that its transcriptional profile is more divergent than would be expected given its experimental condition, as had been hinted previously (Figure 3.3B). (E) PCA plot of gene expression in GSE138058\cite{somasekharan_g3bp1-linked_2020} samples (\gls{pc3} cells). Samples are coloured in accordance to \gls{pl} target (empty vector (green), G3BP1 (red), polysomes (yellow), whole-cell (blue)).  Red and green samples appear intermingled. In this case, however, this pattern is not due to sample mislabelling but rather to another strong transcriptomic effect (visible in the yellow samples on the right side of the plot), which dominates the variance and partially masks the differences between these groups. (F) PCA plot of gene expression in GSE242766\cite{di_timoteo_m6a_2024} samples (SK-N-BE cells). Samples are coloured in accordance to sample type (if the transcriptomes are from \gls{sg} or \gls{wc}, lighter \gls{vs} darker shades, respectively), and to \gls{m6a} status (\gls{ko} (red) and \gls{wt} (blue)). Two major axes of variance are apparent, forming four well-defined quadrants. These axes correspond to sample type (SG vs WC) and m6A status (WT vs KO). In all PCA plots, samples as circles, with the corresponding legend on the plot.}
\end{figure}

\footnotetext{In this case, the discrepancy was indeed caused by sample mislabelling. Visual inspection of the genomic knockout region using \gls{igv}, similar to the approach that will be described in more detail in Chapter 8, revealed the presence of reads within a region that should have been deleted in the \gls{ko} samples. This analysis was communicated to the original authors, who subsequently released a corrigendum\cite{matheny_corrigendum_2023}, although unfortunately without acknowledging this contribution.}

\subsection{Differential Expression Analysis – Linear Modelling}

To formally evaluate gene expression differences between conditions, a useful methodology is to define linear models to describe how experimental variables influence gene expression levels. Consider the example of GSE242766\cite{di_timoteo_m6a_2024} (Figure 3.4F). In this dataset, two main variables are present, which also correspond to the axes of variance observed in the PCA: Sample Type (\gls{wc} or \gls{sg}) and \gls{m6a} state (\gls{wt} or \gls{ko}). A third variable of interest, the presence of an \gls{als}-related mutation, will be ignored here for simplicity.
\\
\\
A direct comparison between two groups, for example SG m6A KO samples and WC m6A WT samples, would confound those variables. Any observed gene expression differences could therefore arise from sample type, \gls{m6a} status, or a combination of both. Linear modelling allows these effects to be partitioned, enabling the contribution of each factor to be estimated independently.
\\
\linebreak
\noindent
To achieve this, gene expression was modelled as a linear function of two explanatory variables: \texttt{SampleType}, encoding the experimental condition, and \texttt{m6AState}, encoding the m6A status. Rather than estimating expression relative to a reference condition, the model was specified without an intercept, such that each level of \texttt{SampleType} receives its own baseline expression estimate. This is equivalent to fitting a separate mean for each group, which facilitates direct pairwise contrasts between conditions without requiring the choice of an arbitrary reference level. In R notation, this corresponds to: 
\begin{verbatim}
	~ 0 + SampleType + m6AState
\end{verbatim}
where \texttt{0} suppresses the intercept, \texttt{SampleType} captures condition-specific mean expression, and \texttt{m6AState} accounts for the effect of \gls{m6a} state across all conditions.
In this formulation, gene expression is explained by the effects of \textit{SampleType} and \textit{m6AState}, allowing their contributions to be estimated separately. This distinction is important for correct biological interpretation. In this dataset, a direct comparison between SG m6A KO and SG m6A WT samples reveals reduced enrichment of longer transcripts in the \gls{ko} samples, as described by the original authors\cite{di_timoteo_m6a_2024}. At first glance, this may suggest that m6A status, rather than gene length, drives \gls{sg} enrichment. However, performing the same comparison in whole-cell samples (WC m6A KO vs WC m6A WT) reveals the same trend. Thus, the \gls{m6a} perturbation consistently produces a transcriptome enriched for shorter \glspl{rna} regardless of sample type. In this context, the \gls{sg} transcriptome likely reflects the underlying \gls{wc} transcriptome, meaning that the apparent reduction in length bias within \glspl{sg} may simply mirror the cellular background (Figure 3.5). 
\\
\linebreak
\noindent
Linear models can also include interaction terms, which test whether the effect of one variable depends on the state of another. In other words, interactions capture deviations from purely additive effects.
\\
\\
As an example, consider GSE138058\cite{somasekharan_g3bp1-linked_2020}. This study used \gls{pl} to identify \glspl{rna} associated with stress granules by targeting \gls{g3bp1} in \gls{pc3} cells under both stressed and unstressed conditions. Additional constructs targeting polysomes, the cytoplasm, and an empty-vector control were also included, each profiled in parallel under both conditions. This experimental design provides appropriate matched controls for each cellular compartment and targeting construct, allowing stress-dependent effects to be distinguished from background labelling or target-specific transcript enrichment.
\\
\linebreak
\noindent
For the analysis of this dataset, gene expression was modelled as a linear function of three components: \texttt{Target}, \texttt{StressState}, and their interaction. As before, the model was specified without an intercept so that each combination of target and stress condition receives its own baseline expression estimate, facilitating direct contrasts between groups. In R notation: 
\begin{verbatim}
	~ 0 + Target * StressState
\end{verbatim}
\noindent
where \texttt{0} suppresses the intercept, and \texttt{Target * StressState} expands to include both main effects and their interaction. That is, the degree to which the effect of stress on gene expression differs depending on which target was used for proximity labelling.
Here, \texttt{Target} represents the proximity-labelling bait (e.g., \gls{g3bp1}, polysome, cytoplasm, or control), and \texttt{StressState} indicates whether cells were stressed or unstressed. Targeting \gls{g3bp1} in unstressed cells captures the baseline \gls{g3bp1} \gls{rna} interactome, whereas targeting \gls{g3bp1} during stress additionally captures \gls{sg}-associated \glspl{rna}. However, stress itself induces widespread transcriptomic changes, meaning that a direct comparison between stressed and unstressed \gls{g3bp1} samples would conflate \gls{sg}-specific interactions with general stress responses. The interaction term addresses this by isolating transcripts whose enrichment in the \gls{g3bp1} interactome is specifically amplified under stress, over and above what stress alone would predict, providing a more precise estimate of the \gls{sg}-associated transcriptome.
\\
\\
In some situations it is also useful to apply variable centering. Centering involves subtracting the mean of a variable, computed across its original defined levels rather than across samples, from each of those levels, so that the variable is expressed relative to its average value. Critically, when interaction terms are included in a linear model, centering alters the interpretation of the main effects: without centering, main effects are estimated at a value of zero for the other variable, which may be outside the range of the data and therefore uninterpretable; with centering, main effects are estimated at the average value of the other variable, which is typically more meaningful\cite{hofmann_centering_1998}. Centering also reduces the numerical collinearity that arises between main effect and interaction terms when variables are not centred, which can otherwise destabilise coefficient estimation\cite{hofmann_centering_1998}. As such, variable centring was applied whenever a linear model with interaction terms was used.
\\
\linebreak
\noindent
Following model fitting, empirical Bayes moderation\cite{smyth_linear_nodate} was applied using the \texttt{eBayes} function from the R package ``limma"\cite{ritchie_limma_2015}. This procedure stabilises gene-wise variance estimates by borrowing information across all genes, producing moderated statistics that are more reliable when the number of replicates is limited\cite{smyth_linear_nodate}. In addition, the empirical Bayes step computes the B-statistic, representing the log-odds that a gene is truly differentially expressed given the observed data, the fitted model, and two prior assumptions: the expected proportion of \glspl{deg} in the dataset, and the minimum log-\gls{fc} required for a gene to be considered differentially expressed\cite{smyth_linear_nodate}.
\\
\linebreak
\noindent
Finally, multiple testing correction must be applied. In transcriptomic analyses, thousands of genes are tested simultaneously for differential expression, meaning that a substantial number of false positives would arise by chance if standard significance thresholds were used.
\\
\\
To control for this issue, the \gls{bh} was used to adjust p-values\cite{benjamini_controlling_1995}. The \gls{bh} controls the \gls{fdr}, the expected proportion of false positives among the results declared significant\cite{benjamini_controlling_1995}. This approach is widely used in \gls{rnaseq} analysis because it balances statistical rigour with sufficient sensitivity to detect biologically meaningful changes in gene expression. 

\subsection{Differential Expression Analysis – Volcano Plots and GSEA}

Volcano plots provide an intuitive representation of differential gene expression analysis (DGEA) because they simultaneously display magnitude and direction of effect, and statistical significance. The horizontal axis reflects the magnitude and direction of expression change, typically log \gls{fc}, while the vertical axis represents the statistical significance of that change. This enables rapid identification of genes combining large-magnitude expression changes with strong statistical significance.
\\
\\
In this analysis, statistical evidence is represented using the B-statistic, which corresponds to the log-odds that a gene is truly differentially expressed given the observed data and the fitted model. Genes appearing toward the upper corners of the plot therefore represent those with the strongest support for differential expression (Figure 3.5A).
\\
\\
To complement gene-level differential expression analysis, we also performed \gls{gsea}\cite{subramanian_gene_2005}, using the fgsea (v1.20.0)\cite{korotkevich_fast_2016} R package (Figure 3.5B). This method evaluates whether groups of biologically related genes exhibit coordinated expression changes, even when individual genes do not reach conventional significance thresholds.
\\
\\
For this analysis, genes were ranked according to their t-statistics obtained from the differential expression model. The t-statistic quantifies the magnitude of expression difference between groups relative to the overall variability observed across samples (i.e., a gene has a high t-statistic when the difference between conditions is large relative to the spread of expression values within those conditions). The ranked gene list was then tested against the Hallmark gene sets, a curated collection designed to reduce redundancy and improve biological interpretability by consolidating overlapping pathways into coherent transcriptional signatures\cite{liberzon_molecular_2015}.
\\
\\
This contrasts with pathway collections such as \gls{kegg}\cite{kanehisa_kegg_2000}, Reactome\cite{croft_reactome_2011}, or Gene Ontology Biological Process\cite{ashburner_gene_2000}, which are typically based on known biochemical or signalling pathway components. While valuable for functional annotation, these pathways often include genes that are mechanistically involved in a process but are not transcriptionally regulated upon pathway activation, limiting their ability to capture coordinated expression changes.
\\
\\
The enrichment analysis produces a \gls{nes}, which quantifies the extent to which genes belonging to a pathway are concentrated near the extremes of the ranked gene list. In the case of t-statistic, the metric used throughout this work, positive \gls{nes} values indicate enrichment among upregulated genes, whereas negative \gls{nes} values indicate enrichment among downregulated genes. By analysing coordinated shifts across gene sets rather than individual genes alone, \gls{gsea} can reveal biologically meaningful patterns that may not be apparent from single-gene differential expression analyses.

\begin{figure}[H]
	\centering
	\includegraphics[width=1\linewidth]{"images/Paper Images/TESE_Fig4.5"}
	\caption[DGEA and GSEA]{\textbf{DGEA and GSEA.} (A) Volcano plots showing \gls{dgea} in GSE242766\cite{di_timoteo_m6a_2024} comparing \gls{sg} samples with their \gls{wc} counterparts in \gls{m6a} \gls{wt} (left) and \gls{m6a} \gls{ko} (right) conditions. Genes are represented as circles. Upregulated transcripts (log2 FC $>$ 2.5 and B-statistic $>$ 10) are highlighted in red, whereas downregulated transcripts (log2 FC $<$ -4 and B-statistic $>$ 10) are shown in blue. Using these thresholds, the same set of genes is highlighted in both conditions. (B) Barplots showing \gls{nes} for Hallmark gene sets comparing \gls{sg} and \gls{wc} samples in \gls{m6a} \gls{wt} (left) and \gls{m6a} \gls{ko} (right) conditions. Under these parameters, the same pathways appear enriched or depleted in both comparisons. Together, panels A and B indicate that the two contrasts produce highly similar results, supporting the conclusion that \gls{m6a} modification does not play a major role in stress granule transcript sequestration.}
	\label{fig:DGEA_GSEA}
\end{figure}

\section{Results and Discussion}

\subsection{Different SG Purification Techniques Yield Different SG Transcriptomes}

As abovementioned, we first compiled publicly available transcriptomic datasets of human \gls{sg} isolated through various methods (Table 3.1).

\begin{table}[H]
	\centering
	\caption[GEO datasets used for SG transcriptome analyses]{\textbf{\gls{geo} datasets used for SG transcriptome analyses}. For each study, the table lists the \gls{gseid}, \gls{sg} isolation method, stress stimulus, target used for \gls{sg} isolation, cell types involved, and associated transcriptome alterations (including transient changes). G3BP refers to G3BP1 alone or both G3BP1 and G3BP2. In the PIC, PHOTON, and FANCI methods, ``target" refers to the antibody used to mark SGs via immunofluorescence, enabling photo-based or fluorescence-activated isolation. DC* – Differential centrifugation without immunoprecipitation; PL** – Proximity labelling-derived technique.}
	\scriptsize 
% inside your table environment, before \begin{tabular}:
	\newlength{\tablewidth}
	\setlength{\tablewidth}{\textwidth - 14\tabcolsep - 8\arrayrulewidth}
	
	\begin{tabular}{|p{0.16\tablewidth}|p{0.10\tablewidth}|p{0.16\tablewidth}|p{0.10\tablewidth}|p{0.10\tablewidth}|p{0.18\tablewidth}|p{0.16\tablewidth}|}
			\hline
			\textbf{Study} & \textbf{Method} & \textbf{Stress} & \textbf{Target} & \textbf{Cell} & \textbf{Transcriptome} & \textbf{GSEID} \\
			\hline
			Khong 2017 & DC & Oxidative & G3BP & U2-OS & G3BP-GFP & GSE99304 \\
			\hline
			Matheny 2019 & DC* & Oxidative & None & U2-OS & Wild-Type & GSE138988 \\
			\hline
			Padrón 2019 & PL & Heat-shock and Hippuristanol & eIF4A & HEK293T & APEX2-eIF4A & GSE121575 \\
			\hline
			Somasekharan 2020 & PL & Oxidative & G3BP & PC-3 & APEX2-G3BP & GSE138058 \\
			\hline
			Matheny 2021 & DC & Oxidative and Hyperosmotic & PABPC1 and G3BP & U2-OS & Wild-Type, G3BP-GFP and G3BP KO & GSE119977 \\
			\hline
			Honda 2021 & PIC & Oxidative & G3BP & HeLa & Wild-Type & GSE143413 \\
			\hline
			Shen 2022 & PL & Oxidative & DDX3X and DDX3Y & HeLa & APEX2-DDX3X and APEX2-DDX3Y & GSE171792 \\
			\hline
			Iadevaia 2022 & DC & Viral Infection & G3BP & U2-OS & G3BP-GFP & GSE171655 \\
			\hline
			Sato 2023 & DC and eCLIP-Seq & Oxidative & G3BP & SH-SY-5Y & Wild-Type & GSE188397 \\
			\hline
			Curdy 2023 & DC & T-Cell Activation & G3BP & Donor T-Cells & Wild-Type & GSE197001 \\
			\hline
			Ren 2023 & PL** & Oxidative and Hyperosmotic & G3BP & HEK293T and U2-OS & G3BP-miniSOG & GSE223295 \\
			\hline
			Xiao 2024 & ARTR-Seq & Oxidative & G3BP & HeLa & Wild-Type & GSE226161 \\
			\hline
			Zhou 2024 & FANCI & Oxidative, Hyperosmotic, Heat-shock and UV Irradiation & G3BP & HeLa & DHX9-eGFP and G3BP-mCherry & GSE218180 \\
			\hline
			Kudrin 2024 & DC & Oxidative & G3BP & HeLa & Wild-Type and NAT10 KO & GSE212380 \\
			\hline
			Di Timoteo 2024 & DC & Oxidative & G3BP & SK-N-BE & Wild-Type, METTL3 KD, and FUS-P525L & GSE242766 \\
			\hline
			Rajachandran 2024 & PHOTON & Oxidative & eIF3eta & HeLa & Wild-Type and METTL3 KO & PRJNA1152621 \\
			\hline
    \end{tabular}
\label{tab:datasets}
\end{table}


\noindent
For each dataset, we compared transcript expression between \gls{sg} and the corresponding whole-cell or cytoplasmic fraction (here collectively referred to as \gls{wc}). \glspl{sg} isolated by \gls{dc} consistently contained longer transcripts on average (Figure 3.6A, Left; Khong \gls{etal} 2017\cite{khong_stress_2017} study shown as example), regardless of stress stimuli or cell type (Figure 3.6B). Gene expression profiles of \gls{dc}-purified \glspl{sg} also exhibited consistently high correlation with each other (Figure 3.6B).
\\
\\
In contrast, \gls{pl}-based datasets showed a weaker opposite \gls{rna} length bias in \glspl{sg} (Figure 3.6A, Right; Somasekharan \gls{etal} 2020\cite{somasekharan_g3bp1-linked_2020} study shown as example) and exhibited minimal correlation between \gls{sg} gene expression profiles, even under the same stress stimulus (Figure 3.6B).
\\
\\
To examine functional enrichment, we performed \gls{gsea} in \glspl{sg} induced by oxidative stress in the \gls{dc} and \gls{pl} example datasets, compared with their respective \gls{wc} fractions (Figure~3.6A,B). In \gls{dc}-purified transcriptomes, gene set enrichment in \glspl{sg} was strongly confounded by transcript length (Fisher's $p < 1.5 \times 10^{-5}$, Odds Ratio $= 23.77$), with several enriched and depleted Hallmarks associated with long (e.g., mitotic spindle) and short (e.g., oxidative phosphorylation) transcripts, respectively (Figure~3.6C). In PL-purified transcriptomes, the observed enrichment of short transcripts in SGs (Figure~3.6A, right) was similarly reflected in Hallmark enrichment analysis (Figure~3.6C; Fisher's $p < 0.0001$, Odds Ratio $= 0.06$).

\begin{figure}[H]
	\centering
	\includegraphics[width=1\linewidth]{"images/Paper Images/Figure1"}
	\caption[SGs purified by different methods have distinct trancriptomes]{}
	\label{fig:Figure1}
\end{figure}

\clearpage

\begin{figure}[t]
	\captionsetup{labelformat=adja-page}
	\ContinuedFloat
	\caption[]{\textbf{SGs purified by different methods have distinct trancriptomes.} (A) Scatter plots of t-statistics of differential expression between SGs and WCs (Y-Axis) \gls{vs} length of gene’s representative \gls{rna} transcript in \gls{nt} (X-Axis, log scale), for \gls{sg} purification by DC\cite{khong_stress_2017} (Left) and PL\cite{somasekharan_g3bp1-linked_2020} (Right); genes as circles, density contour lines in dark red, Spearman's correlation coefficient$\rho$ and associated test's p-value on top left corner. (B) Heatmap of Spearman's correlations between t-statistics of DGEA between SGs and WCs, from different studies, and \gls{rna} length. Stress stimulus, cell type, transcriptome of origin, \gls{sg} purification method, and target molecule used for purification of each study are given below the heatmap. G3BP refers to \gls{g3bp1} alone or both \gls{g3bp1} and \gls{g3bp2}. \gls{mettl3} knockdown and knockout are marked with the same colour. (C) Scatter plot of \gls{gsea} \gls{nes} of Hallmarks between \glspl{sg} and \glspl{wc}, based on genes ranked by t-statistics of differential expression, for \gls{sg} purification by DC\cite{khong_stress_2017} or PL\cite{somasekharan_g3bp1-linked_2020} (SG NES - Y-Axis) \gls{vs} genes ranked by length of representative transcript (RNA Length NES - X-Axis).}
\end{figure}
\noindent
Notably, these datasets differ in cell lines and arsenite concentrations, factors previously suggested to influence \gls{sg} transcript composition. To evaluate whether these biological variables could explain the observed differences, we examined \glspl{wc} from the corresponding cell lines under oxidative stress. These samples showed a correlated transcriptomic stress response (Figure 3.7A-Left), indicating that cell line alone is unlikely to account for the differences observed in SG transcriptomes.
\\
\\
Consistent with this, arsenite-treated U2-OS and PC-3 cell lines share highly upregulated genes, including \textit{\gls{mt1x}} and\textit{ \gls{mt2a}} (metallothioneins protective against oxidative stress), \textit{\gls{dusp1}} (upregulated during oxidative stress), and \textit{\gls{dnajb1}} (mitigates oxidative stress-induced cytotoxicity). Furthermore, \gls{gsea} on \glspl{dgea} between stressed and unstressed states in both cell lines revealed similar Hallmark enrichment (Figure~3.7A-Right; Fisher's $p < 3.5 \times 10^{-4}$, Odds Ratio $= 21.66$). Together, these results suggest that the observed variation in these \gls{sg} transcriptomes is most likely driven by differences in isolation method rather than underlying biological differences between cell lines or stress conditions. While the conserved nature of \glspl{sg} suggests that some degree of consistency across stressors and cell types should be expected, we sought to verify this assumption and to explore alternative explanations that might account for the observed patterns.

\begin{figure}[H]
	\centering
	\includegraphics[width=1\linewidth]{"images/Paper Images/SupFigure1_comp"}
	\caption[U2-OS and PC-3 cells have a similar stress response]{\textbf{U2-OS and PC-3 cells have a similar stress responses.} (A) (Left) Scatter plot of t-statistics of differential gene expression between \glspl{sg} and \glspl{wc} in PC-3\cite{somasekharan_g3bp1-linked_2020} (Y-Axis) and U2-OS\cite{khong_stress_2017} (X-Axis) cell lines exposed to oxidative stress through arsenite;  genes as circles, density contour lines in dark red, Spearman's correlation coefficient $\rho$ and associated test's p-value on top left corner; strongly upregulated genes in both cell lines (t-statistic $>$ 25) are labelled and highlighted in dark red. The dashed line corresponds to the least-squares linear regression fit between the two variables. (Right) Scatter plot of GSEA NES of Hallmarks, based on genes ranked by t-statistics of differential gene expression, in PC-3\cite{somasekharan_g3bp1-linked_2020} (Y-axis)  \gls{vs} U2-OS\cite{khong_stress_2017}  (X-axis) cell lines, exposed to oxidative stress through arsenite. (B) Scatter plots of gene expression (in log2 \gls{cpm}) in \glspl{sg} (Y-Axis) \gls{vs} \glspl{wc} (X-Axis), for \gls{sg} purification by DC\cite{khong_stress_2017}  (left) and PL\cite{somasekharan_g3bp1-linked_2020} (right); genes as circles, density contour lines in dark red, Spearman's correlation coefficient $\rho$ and associated test's p-value on top left corner.}
	\label{fig:Figure2}
\end{figure}

\subsection{Sampling Bias Does Not Explain the RNA Length Bias in DC-derived SGs}

If stress granules strongly and selectively recruited a distinct set of RNAs, their RNA composition might be expected to differ substantially from the whole-cell RNA composition. However, since the two profiles are highly correlated (Figure 3.7B), the stress granule RNA population may largely reflect the general cellular RNA pool. Therefore, the observed enrichment of longer RNAs in DC-purified stress granules may not necessarily reflect a biological property of stress granules themselves, but instead, a technical artefact.
\\
\\
In particular, since \glspl{sg} contain only about 10\% of total cellular \gls{mrna}\cite{khong_stress_2017}, the \glspl{rna} detected in \glspl{sg} may represent only a partial sampling of the whole-cell transcriptome. In such smaller populations, random fluctuations can exaggerate the representation of certain transcript features, such as \gls{rna} length, a phenomenon conceptually related to sampling bias\cite{phipson_gene_2017}.
\\
\\
To evaluate whether technical or sampling effects could account for the observed length bias, we performed \textit{in silico} \gls{rnaseq} simulations using empirical \gls{rna} expression levels and fragment size distributions from the original \gls{sg} and \gls{wc} samples (schematic shown in Figure 3.8; see section 3.1.3). These simulations allowed us to test whether sampling or \gls{rnaseq} processes alone could reproduce the observed enrichment patterns.
\\
\begin{figure}[H]
	\centering
	\includegraphics[width=1\linewidth]{"images/Paper Images/TESE_Fig4.8"}
	\caption[Mimi-Seq Simulation Pipeline]{\textbf{Mimi-Seq simulation pipeline.} The diagram depicts the experimental workflow from the original study for obtaining WC and SG samples\cite{khong_stress_2017} (steps in grey) and the additional \textit{in silico} steps performed in our analyses (steps in black). }
	\label{fig:Figure4.8}
\end{figure}
\noindent
Two potential sources of bias were considered. First, small \gls{rna} populations can be subjected to technical sampling bias, where random over- or under-representation of transcripts can distort apparent transcript distributions. Second, \gls{rnaseq} library preparation involves fragmentation, which generates more fragments, and therefore more reads, for longer transcripts, an effect that standard normalisation methods that account for transcript length do not fully correct\cite{zhao_tpm_2021, mandelboum_recurrent_2019}.
\\
\\
As control, we simulated \gls{rnaseq} using \gls{wc} and \gls{sg} expression data with their respective fragmentation profiles (i.e., how specific transcripts fragment in each sample type), reproducing empirical fragment length distributions(\gls{vi} section 3.1.3\cite{khong_stress_2017}) (Figure 3.9A). Additional simulations in which expression data from one population were paired with the alternate fragmentation profile demonstrated that fragmentation patterns depended only on the input expression, confirming equivalence of \gls{sg} and \gls{wc} fragmentation profiles (Figure 3.9B).
\\
\\
Finally, we simulated \gls{wc} and \gls{sg} samples by drawing reads from the full \gls{wc} transcriptome, with \gls{sg} samples representing a 10\% subset of that expression. These simulated \gls{sg} transcriptomes exhibited a negative length bias, in contrast to the positive bias in \gls{dc}-purified \gls{sg} transcriptomes (Figure 3.9C). This reversal is most likely due to the high abundance of short\textit{ RN7SL1–3} \glspl{rna} in the original \gls{wc} transcriptome, which constitute approximately 50\% of total reads\cite{khong_stress_2017}, and outweigh contributions from longer transcripts in the simulation.
\\
\\
Even in this simplified scenario, designed to maximise technical enrichment of longer \glspl{rna}, the positive length bias in \glspl{sg} was not reproduced. In fact, sequencing depth and coverage of \glspl{sg} (as explained in section 3.2.2) in the original study\cite{khong_stress_2017} appear sufficient to prevent this phenomenon (Figure 3.9D).
\\
\begin{figure}[H]
	\centering
	\includegraphics[width=1\linewidth]{"images/Paper Images/TESE_Fig4.9"}
	\caption[Transcriptomes of DC-isolated SGs show no evidence for RNA fragmentation-associated sampling bias]{\textbf{Transcriptomes of DC-isolated SGs show no evidence for RNA fragmentation-associated sampling bias.} (A) Relative distribution (density) of fragment length (\gls{nt}) in original samples\cite{khong_stress_2017} \gls{rnaseq} (Left) and simulated samples (Right). (B) Relative distribution (density) of representative \gls{rna} fragment length (\gls{nt}) in simulated samples, using \gls{wc} (Left) or \gls{sg} (Right) expression as basis, and \gls{wc} or \gls{sg} fragmentation distributions as basis. (C) Scatter plot of t-statistics of \gls{dgea} between \gls{sg} and \gls{wc} simulated samples (Y-Axis) and representative \gls{rna} transcript length (\gls{nt}) (X-Axis, log scale); genes as circles, density contour lines in dark red, Spearman's correlation coefficient $\rho$ and associated test's p-value are also shown. The original WC expression\cite{khong_stress_2017} was used as the basis for both simulated \gls{wc} and \gls{sg} samples, with their respective fragment size distribution as the basis. (D) Cumulative distribution of the percentage of \gls{rnaseq} reads mapping to the increasing number of genes, ordered by decreasing read coverage, in the original dataset\cite{khong_stress_2017}.  More genes account for 50\% of the reads in \gls{sg} samples than their \gls{wc} counterparts, suggesting greater \gls{rna} library complexity.}
	\label{fig:Figure4.9}
\end{figure}
\noindent
Collectively, these results suggest that the positive length bias in \gls{dc}-purified \gls{sg} transcriptomes is unlikely to arise from either chance sampling or \gls{rnaseq} technical effects owing to the lower number of \gls{rna} molecules on \gls{sg}, when compared to the remainder of the cell.

\subsection{Simulation of the Centrifugation Process Yields RNA Length Bias Compatible with that Observed in DC-isolated SGs}

\glspl{rna} can naturally form condensates in the cell independently of their localisation to \glspl{sg}\cite{glauninger_stressful_2022, glauninger_transcriptome-wide_2025, van_treeck_rna_2018}\footnote{Perhaps the strongest support for this idea comes from experiments by the same research group that first isolated stress granules, in which cell-free RNA was shown to spontaneously assemble into \gls{sg}-like condensates that share a similar transcriptomic composition\cite{van_treeck_rna_2018}. These observations were largely interpreted as evidence that the properties governing \gls{sg} assembly are ``encoded" at the \gls{rna} level. However, an alternative interpretation is also possible: that the intrinsic propensity of \glspl{rna} to self-assemble may generate condensates with \gls{sg}-like transcriptomes even outside physiological stress conditions. Because these assemblies occur in cell-free, stress-free environments (albeit under molecular crowding conditions) it is conceivable that similar \gls{rna} self-assembly could also occur during differential centrifugation. In this scenario, the observed transcript composition would not necessarily reflect \gls{sg}-specific sequestration but rather the intrinsic aggregation properties of \gls{rna} molecules themselves.}. Thus, \gls{dc}, which isolates condensate-rich fractions prior to immunoprecipitation, may capture \gls{rna} condensates other than \glspl{sg}. To investigate this possibility, we simulated centrifugation in silico using \gls{wc} expression data\cite{khong_stress_2017}, modelling \gls{rna} precipitation as primarily dependent on transcript molecular weight under the assumption of uniform condensate shape (Figure 3.10A). Although \glspl{rna} sediment as part of condensates, we modelled each transcript as sedimenting independently, using molecular weight as a proxy for condensate behaviour (see section 3.1.4). This approach generates a gene expression profile attributable solely to the technical effects of centrifugation, which we interpret as putative \glspl{sg}.
\\
\\
The simulated \glspl{sg} displayed a positive \gls{rna} length bias relative to \gls{wc} samples, closely matching that observed in the original \gls{dc} data (Figure~3.10B). \gls{gsea} was similarly confounded by transcript length (Figure~3.10C; $p$-value $< 1 \times 10^{-9}$, Odds Ratio $= 44.18$), consistent with earlier observations (Figure~3.7A-Right).
\\

\begin{figure}[H]
	\centering
	\includegraphics[width=1\linewidth]{"images/Paper Images/Figure3"}
	\caption[Length bias in DC-isolated SGs is compatible with inherent RNA condensation]{}
	\label{fig:Figure4.10}
\end{figure}

\clearpage


\begin{figure}[t]
	\captionsetup{labelformat=adja-page}
	\ContinuedFloat
	\caption[]{\textbf{Length bias in DC-isolated SGs is compatible with inherent RNA condensation.} (A) Schematic representation of standard \gls{dc} protocol alongside our \textit{in silico} pipeline for simulating the centrifugation process. d = distance travelled by each RNA during centrifugation; Mw = Molecular weight of each \gls{rna} transcript, in g/mol; g = centrifugation acceleration, in g; t = time of centrifugation, in seconds.  (B) Scatter plot of t-statistics of differential gene expression between simulated \glspl{sg}, and \glspl{wc} (Y-Axis) \gls{vs} length of gene's representative \gls{rna} transcript (\gls{nt}) (X-Axis, log scale). (C) Scatter plot of \gls{gsea} \gls{nes} of Hallmark gene sets between simulated \gls{sg} and \gls{wc} samples, with genes ranked by t-statistics of differential expression (Y-Axis) \gls{vs} genes ranked by length (\gls{nt}) (X-Axis). Fisher’s test \gls{or} and associated p-value on top right corner. (D) Scatter plots of t-statistics of differential gene expression in simulated \glspl{sg} (Y-Axis) and \gls{dc}-purified \glspl{sg} (X-Axis), relative to \gls{wc} samples; genes as circles coloured by RNA length in \gls{nt}; density contour lines in dark red, Spearman's correlation coefficient $\rho$ and associated test's p-value on top left corner. (E) Scatter plots of t-statistics of differential gene expression in \gls{dc}-purified \glspl{sg} \gls{vs} simulated \glspl{sg} (Y-Axis) \gls{vs} length of gene's representative \gls{rna} transcript (\gls{nt}) (X-Axis, log scale); transcripts as circles; density contour lines in dark red, Spearman’s correlation coefficient $\rho$ and associated test’s p-value on top left corner. (F) Scatter plots of t-statistics of differential gene expression between stressed and unstressed \gls{rnp} samples (Y-Axis), and between stressed and unstressed \gls{wc} samples (X-Axis). Genes as circles; density contour lines in dark red, Spearman's correlation coefficient $\rho$ and associated test's p-value on top left corner (B, D, E, F). The dashed line corresponds to the least-squares linear regression fit between the two variables (D, E, F). The original \gls{wc} expression from GSE99304\cite{khong_stress_2017} was used as basis for simulated \gls{sg} samples.}
\end{figure}
\noindent
To assess the similarity between simulated and \gls{dc}-purified \glspl{sg}, we compared their differentially expressed genes relative to \gls{wc} samples and observed a strong correlation (Figure 3.10D)\footnote{Even though the correlation is considerable, it is far from perfect, as the model itself relies on very rough approximations (e.g. every RNA transcript behaving as a sphere). This imperfection, however, can be seen as an advantage. Even a deeply flawed formulation can recapitulate the DC results, to a certain extent. Naturally, this does not prove that the signal is solely methodological, but does warrant investigation of this possibility.}. Because the simulated \glspl{sg} were generated solely from \gls{wc} expression data and lack any \gls{sg}-specific biological signal, they represent the centrifugation-driven component of the \gls{dc} method. Comparing these simulated \glspl{sg} with the experimentally derived SG fractions allows identification of transcripts enriched beyond what would be expected based on molecular weight alone. After this correction, the resulting SG-enriched RNA profiles showed no detectable RNA length bias (Figure 3.10E). This suggests that the length bias reported in previous studies likely arises from centrifugation-related artefacts during \gls{sg} purification. For simplicity, we refer to these corrected SG RNA profiles as ``SG transcriptomes” throughout the remainder of this work.
\\
\\
It is important to acknowledge that this centrifugation simulation relies on coarse approximations and may not fully capture the true behaviour of \gls{rna} during centrifugation. This limitation may explain why the transcripts identified as differentially enriched often correspond to extreme transcript lengths, either very short or very long (Figure 3.10E, top left and bottom right).
\\
\\
If the \gls{sg} signal is indeed confounded by natural \gls{rna} condensation, a more appropriate control for \gls{dc}-purified \glspl{sg} requires samples processed identically but lacking \glspl{sg}. In Matheny \gls{etal} 2019\cite{matheny_transcriptome-wide_2019}, \gls{wc} and \gls{rnp} granule fractions were collected under both stressed and unstressed conditions using \gls{dc}, but without immunopurification. In this dataset, stressed \gls{rnp} fractions contain \glspl{sg} together with other \gls{rna} granules such as \glspl{pb}, which are also present in unstressed cells and are associated with \gls{rna} degradation. In contrast, unstressed \gls{rnp} fractions lack \glspl{sg}.
\\
\\
If \gls{sg}-specific \gls{rna} sequestration occurs beyond the effects of stress and centrifugation, transcriptomic changes between stressed and unstressed \gls{rnp} fractions would be expected to deviate from those observed in the corresponding \gls{wc} comparisons, with transcripts specifically recruited to or excluded from \gls{sg}s shifting the \gls{sg} comparison relative to the \gls{wc} baseline. However, the two comparisons are strongly correlated (Spearman's ρ = 0.71; Figure 3.10F): transcripts upregulated in stressed vs. unstressed WC are similarly upregulated in stressed vs. unstressed SG, and vice versa. This concordance indicates that the SG transcriptome largely mirrors the cellular transcriptome under both conditions, with minimal transcript-specific sequestration. Variance analysis further supports this interpretation, showing that SG-specific effects are smaller than those attributable to stress and centrifugation (Figure 3.11).
\\
\\
Notably, this study did not include immunopurification of a core \gls{sg} protein, which may dilute \gls{sg}-specific signals among other \gls{rna} condensates and contribute to the lack of clear \gls{sg}-specific enrichment. In addition, this protocol includes one fewer centrifugation step than other \gls{dc}-based methods, which may explain the weaker correlation with \gls{rna} length observed in this dataset compared to other \gls{dc} studies using the same cell types and stress conditions (Figure 3.6B).
\begin{figure}[H]
	\centering
	\includegraphics[width=1\linewidth]{"images/Paper Images/Sup3"}
	\caption[Stress granules presence contribution to variance in RNP granules]{\textbf{Stress granules presence contribution to variance in RNP granules.} Relative contributions of different sources to gene expression variance in the GSE138988\cite{matheny_transcriptome-wide_2019} dataset. \gls{sg} contribute less to gene expression variance than the stress itself or the centrifugation effect.}
	\label{fig:Figure4.11}
\end{figure}
\noindent

\subsection{RNA Modifications Have a Limited (Direct) Effect on SG Transcript Enrichment}

Beyond \gls{rna} length and translation efficiency\cite{khong_stress_2017}, \gls{rna} modifications such as \gls{m6a} and \gls{ac4c} have also been associated with \gls{sg} enrichment in datasets obtained through \gls{dc}\cite{fu_m6a-binding_2020, ries_m6a_2023, kudrin_n4-acetylcytidine_2024, di_timoteo_m6a_2024}. To evaluate whether these modifications genuinely influence \gls{sg} transcript enrichment, we compared differential gene expression between \gls{sg} and \gls{wc} fractions in \gls{ac4c}- and \gls{m6a}-\gls{ko} cells relative to their \gls{wt} counterparts (Figure 3.12A).

\begin{figure}[h]
	\centering
	\includegraphics[width=1\linewidth]{"images/Paper Images/Figure4_2"}
	\caption[DC-purified SG-enriched transcripts with RNA modifications are compatible with centrifugation biases]{}
	\label{fig:Figure4.12}
\end{figure}

\clearpage


\begin{figure}[t]
	\captionsetup{labelformat=adja-page}
	\ContinuedFloat
	\caption[]{\textbf{AC4C, but not m6A, is associated with altered SG transcript enrichment, 
			driven by RNA length.}
		(A) Scatter plots of t-statistics of differential gene expression in \gls{wt} \glspl{sg} 
		\gls{vs} \glspl{wc} (Y-axis) and in \glspl{sg} \gls{vs} \glspl{wc} in cells without \gls{rna} 
		modifications (X-axis) -- \gls{ac4c} KO\cite{kudrin_n4-acetylcytidine_2024} (Left) and 
		\gls{m6a} KO\cite{di_timoteo_m6a_2024} cells (Right); transcripts as circles, density 
		contour lines in dark red, Spearman's correlation coefficient $\rho$ and associated test's 
		$p$-value on the top left corner.
		(B) Volcano plots of the transcriptomic effect (given by linear modelling with interaction) 
		of \gls{ac4c} (Left) and \gls{m6a} (Right) modifications in \gls{sg} samples.
		(C) Violin plots of residuals (the difference between the actual value and the value 
		predicted by the linear model applied for any given transcript) as a function of the number 
		of \gls{ac4c} modifications per transcript; Spearman's correlation coefficient $\rho$ and 
		associated test's $p$-value shown on the top left corner.
		(D) Scatter plot of residuals (Y-axis) \gls{vs} transcript length (nt) (X-axis); 
		transcripts as circles, density contour lines in dark red, Spearman's correlation 
		coefficient $\rho$ and associated test's $p$-value shown on the top left corner. The dashed line corresponds to the least-squares linear regression fit between the two variables (A, D).
		For all panels concerning \gls{ac4c}, the GSE212380 dataset was used\cite{kudrin_n4-acetylcytidine_2024}. For 
		those concerning \gls{m6a}, the GSE242766 dataset was 
		used\cite{di_timoteo_m6a_2024}}
\end{figure}
\noindent
Using the residuals from a linear model predicting transcript sequestration within \glspl{sg} as a measure of modification-dependent effects (an interaction model, as described in 3.2.5)\footnote{It is worth noting that in the original studies, the reported effects of \gls{rna} modifications were likely driven in part by differences already present in \gls{wc} samples, arising from direct comparisons between different sample types. An example illustrating this analytical pitfall is shown in Section 3.2.6.}, we assessed whether the absence of each modification systematically shifted \gls{sg} enrichment. For \gls{m6a}, no discernible effect was observed (Figure 3.12B - Right). For \gls{ac4c}, however, a modest but detectable shift was present, suggesting that \gls{ac4c} may influence \gls{sg} transcript enrichment (Figure 3.12B - Left).
\\
\\
Given this finding, we sought to determine what underlies the \gls{ac4c}-associated effect. 
We first examined whether the number of \gls{ac4c} modifications per transcript could account 
for it, as the original authors posit\cite{kudrin_n4-acetylcytidine_2024}, reasoning that 
transcripts bearing more modifications would be most affected if the modification itself were 
causative. As such, we compared these with the residuals (i.e., how much a sample changes 
from what was expected from the dashed line in Figure 3.12A - Left). However, only a very weak correlation was observed between the number of \gls{ac4c} modifications and the magnitude of the enrichment shift (Figure 3.12C). Additional \textit{in silico} simulations taking into account the extra weight from these modifications showed only limited correlation with enrichment (see section 3.1.4).  We subsequently tested \gls{rna} length as an alternative explanatory 
variable, and found that it substantially accounts for the observed differences (Figure 3.12D).
\\
\\
These results present an unexpected finding. \gls{ac4c} appears to be associated with altered \gls{sg} enrichment, yet this effect does not appear to act through \gls{ac4c}-modified transcripts themselves, but rather through other transcripts. Why \gls{rna} length would emerge as a distinguishing variable between \gls{ac4c} \gls{ko} and \gls{wt} conditions, given that both underwent the same differential centrifugation pipeline, remains unclear. Nevertheless, the data suggest that the \gls{ac4c}-associated signal in \gls{sg} enrichment studies reflects a length-dependent bias rather than a direct effect of the modification on \gls{sg} sequestration.

\subsection{Differential Centrifugation Yields Similar Transcriptomes for Stress Granules and other RNP Granules}

Based on the preceding results, we hypothesise that the \gls{sg} transcriptome isolated by \gls{dc} is substantially confounded by the centrifugation process itself rather than reflecting exclusive biological sequestration. Within this hypothesis, other \gls{rna} granules isolated using similar \gls{dc} protocols, regardless of their biological identity, should display comparable transcriptomic profiles, specifically enrichment for longer \glspl{rna}.
\\
\\
In the GSE138988\cite{matheny_transcriptome-wide_2019} dataset, stressed \glspl{pb}, cytoplasmic ribonucleoprotein condensates involved in mRNA decay and translational repression\cite{matheny_transcriptome-wide_2019}, and bulk \gls{rnp} granules were purified by \gls{dc}, with \glspl{pb} further distinguished by an additional immunoprecipitation step using \gls{edc3}. The original GSE99304\cite{khong_stress_2017} study used a similar three-round \gls{dc} protocol to isolate \glspl{sg}, employing G3BP for immunoprecipitation. Despite representing distinct cytoplasmic granules, \gls{dc}-purified \glspl{sg} and \glspl{pb} displayed similar transcriptomes (Figure 3.13A), as also noted by the authors\cite{matheny_transcriptome-wide_2019}. Moreover, both \glspl{sg} and \glspl{pb} were more similar to each other than to bulk \gls{rnp} granules purified using only two rounds of centrifugation (Figure 3.13B–C). It is unclear whether this pattern reflects the presence of multiple condensate types within the bulk \gls{rnp} preparations or, alternatively, the effect of the additional centrifugation step shared by the \gls{sg} and \gls{pb} protocols.
\\
\\
Further support for the centrifugation bias hypothesis comes from comparisons with datasets generated using non-centrifugation-based isolation methods. In GSE138988\cite{matheny_transcriptome-wide_2019}, \gls{dc}-isolated stressed \glspl{pb} showed lower similarity to unstressed \glspl{pb} isolated by \gls{faps} in Hubstenberger \gls{etal} 2017\cite{hubstenberger_p-body_2017} than to \gls{dc}-isolated \glspl{sg}, as pointed by the authors\cite{matheny_transcriptome-wide_2019}. We confirmed this observation using an independent \gls{faps} dataset, Kodali \gls{etal} 2024\cite{kodali_rna_2024} (Figure 3.13D). Although these \gls{pb} datasets remain correlated, stressed \glspl{pb} isolated by \gls{dc} still show higher similarity to \gls{dc}-isolated \glspl{sg} or \gls{rnp} fractions than to \glspl{pb} isolated by \gls{faps} (Figure 3.13A–D).
\\
\\
In Zhou \gls{etal} 2024 (GSE218180)\cite{zhou_rna_2024}, stress granules were isolated using \gls{fanci}, a flow cytometry–based approach conceptually similar to \gls{faps}. Even after batch correction for cell type and stress, \gls{fanci}-isolated \glspl{sg} remained less similar to \gls{dc}-isolated \glspl{sg} than the latter were to \gls{dc}-isolated \glspl{pb}, despite \glspl{sg} and \glspl{bp} being biologically distinct structures (Figure 3.13E).
\\
\\
The authors of the GSE138988\cite{matheny_transcriptome-wide_2019} study proposed that cellular stress increases the pool of non-translating \glspl{mrna}, thereby shifting the \gls{pb} transcriptome toward a more \gls{sg}-like composition. However, our reanalysis does not fully support this explanation. To further investigate this possibility, we integrated the Zhou \gls{etal} 2024 and Kodali \gls{etal} 2024 studies, both generated using flow cytometry–based isolation methods in \gls{hela} cells. Because both studies include unstressed WC controls, batch correction for stress could be applied. After this adjustment, similarities between \glspl{sg} and \glspl{pb} isolated by FAPS/FANCI remained detectable but were substantially weaker than those observed in \gls{dc}-based datasets (Figure 3.13F).
\\

\begin{figure}[H]
	\centering
	\includegraphics[width=1\linewidth]{"images/Paper Images/Figure5"}
	\caption[Samples isolated through DC are highly correlated]{\textbf{Samples isolated through DC are highly correlated.}
		(A) Scatter plot of t-statistics of \gls{dgea} in \glspl{sg} \gls{vs} \glspl{wc}\cite{khong_stress_2017} (Y-Axis) and in \glspl{pb} \gls{vs} \glspl{wc}\cite{matheny_transcriptome-wide_2019} (X-Axis). (B) Scatter plot of t-statistics of \gls{dgea} in \glspl{sg} \gls{vs} \glspl{wc}\cite{matheny_transcriptome-wide_2019} (Y-Axis) and in \gls{rnp} granules \gls{vs} \glspl{wc}\cite{matheny_transcriptome-wide_2019} (X-Axis). (C) Scatter plot of t-statistics of \gls{dgea} in \glspl{pb} \gls{vs} \glspl{wc}\cite{matheny_transcriptome-wide_2019} (Y-Axis) and in \gls{rnp} granules \gls{vs} \glspl{wc}\cite{matheny_transcriptome-wide_2019} (X-Axis). (D) Scatter plot of t-statistics of \gls{dgea} in \glspl{pb} purified by FAPS\cite{kodali_rna_2024} (Y-Axis) and by \gls{dc}\cite{matheny_transcriptome-wide_2019}, compared to \glspl{wc} (X-Axis). (E) Scatter plot of t-statistics of \gls{dgea} in \glspl{sg} purified by \gls{dc}\cite{khong_stress_2017} (Y-Axis) and by FAPS/FANCI\cite{zhou_rna_2024} (X-Axis), comparing to \glspl{wc}, after batch correction for cell type and stress. (F) Scatter plot of t-statistics of \gls{dgea} in \gls{sg} (Y-Axis) and \gls{pb} (X-Axis) purified by FAPS/FANCI\cite{kodali_rna_2024, zhou_rna_2024} comparing to \glspl{wc}, after batch correction for stress; In all panels, genes as circles coloured by \gls{rna} length (nt); density contour lines in dark red, Spearman's correlation coefficient $\rho$ and associated test's p-value on top left corner. The dashed line corresponds to the least-squares linear regression fit between the two variables.}
	\label{fig:Figure4.13}
\end{figure}
\noindent
It is important to note that this batch correction approach assumes that stress affects \glspl{sg} and \glspl{pb} in comparable ways, which may not fully capture more complex biological interactions. Nevertheless, because stressed \gls{pb} datasets isolated by \gls{faps} are currently unavailable, it remains unclear whether the divergence between \gls{dc} and flow cytometry–based datasets reflects genuine stress-induced remodeling of \glspl{pb} or, alternatively, technical artefacts introduced by centrifugation, such as length-dependent \gls{rna} bias.
\\
\\
These batch correction methods are explained in more detail in their own chapter, in section 5.2.2.

\subsection{Revisiting Stress Granule Transcriptomes Reveals Mitochondrial RNA Enrichment Despite Methodological Bias}

To further evaluate the centrifugation hypothesis, we analysed publicly available Sed-seq data, which measures \gls{rna} sedimentation by sequencing pelleted and unpelleted fractions from stressed and unstressed yeast cells\cite{glauninger_transcriptome-wide_2025}. These data allowed us to derive transcript-specific \gls{s} (see section 3.1.4), providing a more accurate estimate of \gls{rna} sedimentation behaviour than our initial approximations based solely on \gls{mw}. Interestingly, mitochondrially encoded transcripts displayed higher sedimentation coefficients than expected based on their length alone (Figure 3.14A), suggesting that these have distinct properties from the nuclear encoded transcripts. In this context, sedimentation coefficients (Svedberg units, S) describe the rate at which RNA-containing particles sediment in a centrifugal field under defined conditions, reflecting their effective size, shape, and density-dependent frictional properties\cite{svedberg_ultracentrifuge_1940}.
\\
\\
Applying these experimentally derived coefficients to simulate \gls{sg} transcript sequestration relative to whole-cell samples in an independent yeast dataset\cite{khong_stress_2017} produced a strong correlation with the observed \gls{sg} transcriptome (Figure 3.14B). These results support the centrifugation hypothesis in an experimental system, although in a different organism.
\\
\\
We next extended this framework to human cells by deriving empirical formulas linking \gls{s} to \gls{rna} length. Separate equations were used for nuclear and mitochondrial transcripts, since mitochondrial \glspl{rna} consistently displayed higher sedimentation coefficients than nuclear transcripts of comparable length (Figure 3.14A), possibly reflecting their polycistronic organisation. These formulas enabled estimation of theoretical \gls{s} for human transcripts and simulation of their expected sedimentation behaviour, producing an \textit{in silico} model of putative \gls{sg} fractions.
\\
\\
Simulations based on \gls{s} more closely matched empirical \gls{sg} data than our earlier weight-based model (Figure 3.14C vs Figure 3.10D). Comparing observed \gls{sg} enrichment with these simulations allowed identification of transcripts enriched or depleted beyond what would be expected from centrifugation alone, thereby defining a corrected \gls{sg} signal that reflects sequestration independent of sedimentation effects. As a proof of concept, this approach recovered known \gls{sg} markers such as \textit{\gls{norad}} as enriched and \textit{\gls{gapdh}} as depleted, consistent with previous \gls{rna}-\gls{fish} validation\cite{khong_stress_2017} (Figure 3.14D). The corrected \gls{sg} signal showed only a modest correlation with \gls{rna} length (Figure 3.14E), substantially weaker than reported in the original study\cite{khong_stress_2017}. \gls{gsea} of the corrected \gls{sg} profile revealed enrichment of mitotic spindle–related transcripts (Figure 3.14F), consistent with the possibility that \glspl{sg} selectively sequester cell-cycle–associated \glspl{rna} during stress. However, the presence of additional signatures, including downregulation of \gls{tnfa} signalling via \gls{nfkb} and upregulation of metabolic pathways, complicates this interpretation and leaves the functional implications of \gls{sg} sequestration unresolved.
\\

\begin{figure}[H]
	\centering
	\includegraphics[width=1\linewidth]{"images/Paper Images/TESE_Fig4.14"}
	\caption[Method-specific corrections for SGs transcriptomes reveal distinct outcomes (Part 1)]{}
	\label{fig:Figure4.14}
\end{figure}

\clearpage


\begin{figure}[t]
	\captionsetup{labelformat=adja-page}
	\ContinuedFloat
	\caption[]{\textbf{Method-specific corrections for SGs transcriptomes reveal distinct outcomes (Part 1).}
		(A) Scatter plot of yeast \gls{s} values obtained from Sed-seq data\cite{glauninger_transcriptome-wide_2025} (Y-Axis) and cubic root of yeast \gls{rna} length (\gls{nt}) (X-Axis); Each circle denotes a transcript, with mitochondrially encoded transcripts coloured in cyan; density contour lines in dark red, Spearman's correlation coefficient $\rho$ and associated test's p-value on top left corner. (B) Scatter plot of t-statistics of differential gene expression of simulated yeast \gls{sg} \gls{vs} \gls{wc} (Y-Axis), with \gls{s} estimations based on Sed-seq data\cite{glauninger_transcriptome-wide_2025}, and t-statistics of differential gene expression in yeast \glspl{sg} purified by \gls{dc}\cite{khong_stress_2017} (X-Axis); transcripts as circles coloured by \gls{rna} length; density contour lines in dark red, Spearman's correlation coefficient $\rho$ and associated test's p-value on top left corner. (C) Scatter plot of t-statistics of differential gene expression of simulated \glspl{sg} \gls{vs} \glspl{wc} (Y-Axis), with \gls{s} estimations based on Sed-seq data\cite{glauninger_transcriptome-wide_2025}, and t-statistics of differential gene expression of \gls{dc}-purified \glspl{sg} \gls{vs} \glspl{wc}\cite{khong_stress_2017} (X-Axis); transcripts as circles coloured by \gls{rna} length; density contour lines in dark red, Spearman's correlation coefficient $\rho$ and associated test's p-value on top left corner. (D) Volcano plot of DGEA (Y-Axis – B-Statistic; X-Axis - log2 Fold-Change) between original \glspl{sg}\cite{khong_stress_2017} and simulated \glspl{sg}, based on Sed-Seq data\cite{glauninger_transcriptome-wide_2025}; transcripts as circles, with commonly used \gls{rna}-\gls{fish} controls\cite{khong_stress_2017} in red. Vertical dashed lines correspond to -0.75 and 0.75 log2 \gls{fc}. Horizontal dashed line corresponds to B-statistic = 0. (E) Scatter plot of \gls{rna} length (X-Axis) and  t-statistics of differential gene expression of Simulated \glspl{sg} \gls{vs} \gls{dc}-Purified \glspl{sg}\cite{khong_stress_2017} (thus what we consider to be “\textit{Centrifugation Corrected \glspl{sg}}") (Y-Axis); transcripts as circles; density contour lines in dark red, Spearman's correlation coefficient $\rho$ and associated test's p-value on top left corner. (F) \gls{gsea} \gls{nes} of Hallmark gene sets, ranked by t-statistic, in Simulated \glspl{sg} \gls{vs} \gls{dc}-Purified \glspl{sg}\cite{khong_stress_2017}. The dashed line corresponds to the least-squares linear regression fit between the two variables (A, B, C, E).}
\end{figure}
\noindent
To further refine \gls{sg} composition in human cells, we applied a correction strategy accounting for both centrifugation and stress effects using the          GSE138988\cite{matheny_transcriptome-wide_2019} dataset. This dataset includes stressed and unstressed \glspl{rnp} and \gls{wc} samples generated using differential centrifugation, allowing the centrifugation component of the protocol to be controlled. However, direct comparison between stressed and unstressed \gls{rnp} fractions is expected to reflect global stress responses rather than \gls{sg}-specific sequestration. To disentangle these effects, we applied linear models accounting for both stress and centrifugation (see section 3.2.5). This analysis produced a refined \gls{sg} profile in which \textit{\gls{norad}} was enriched and \textit{\gls{gapdh}} and \textit{\gls{actb}} were depleted, again consistent with \gls{rna}-\gls{fish} validation\cite{khong_stress_2017} (Figure 3.15A). The resulting \gls{sg} signal showed a positive but substantially weaker correlation with \gls{rna} length than previously reported\cite{matheny_transcriptome-wide_2019} (Figure 3.15B).
\\
\\
\gls{pl} approaches do not rely on centrifugation and therefore avoid sedimentation biases, but introduce other potential confounding factors. For example, \gls{g3bp1} does not localise exclusively to \glspl{sg}, meaning that \gls{pl} approaches may capture \glspl{rna} outside bona fide \glspl{sg}. Although this limitation cannot be directly controlled, we applied linear modelling strategies to account for additional confounding variables (see section 3.2.5). In some \gls{pl}-purified \gls{sg} datasets\cite{somasekharan_g3bp1-linked_2020}, the resulting \gls{sg}-specific signal displayed a negative correlation with \gls{rna} length (Figure 3.15C), and gene set enrichment analysis revealed enrichment of metabolic and stress-response pathways (Figure 3.15D), consistent with selective transcript sequestration (see section 1.1.4).
\\


\begin{figure}[H]
	\centering
	\includegraphics[width=1\linewidth]{"images/Paper Images/TESE_Fig4.15"}
	\caption[Method-specific corrections for SGs transcriptomes reveal distinct outcomes (Part 2)]{}
	\label{fig:Figure4.15}
\end{figure}

\clearpage

\begin{figure}[t]
	\captionsetup{labelformat=adja-page}
	\ContinuedFloat
	\caption[]{\textbf{Method-specific corrections for SGs transcriptomes reveal distinct outcomes (Part 2).} (A) Volcano plot of differential gene expression (Y-Axis – B-Statistic; X-Axis - log2 Fold-Change) of interception of stress with \gls{rnp} sample types, given by linear modelling (i.e., corresponding to \gls{dc}-isolated \gls{rnp} granule stress-specific sequestration\cite{matheny_transcriptome-wide_2019}); transcripts as circles, with commonly used \gls{rna}-\gls{fish} controls\cite{khong_stress_2017} in red. Vertical dashed lines correspond to -0.75 and 0.75 log2 \gls{fc}. Horizontal dashed line corresponds to B-statistic = 0. (B) Scatter plot of t-statistics of differential gene expression in linear modelled \gls{dc}-isolated \gls{rnp} granules stress-specific sequestration\cite{matheny_transcriptome-wide_2019} (Y-Axis) and \gls{rna} length (X-Axis); transcripts as circles; density contour lines in dark red, Spearman's correlation coefficient $\rho$ and associated test's p-value on top left corner. (C) Scatter plot of t-statistics of differential gene expression in linear modelled \gls{pl}-isolated \glspl{sg}\cite{somasekharan_g3bp1-linked_2020} (Y-Axis) and \gls{rna} length (X-Axis); transcripts as circles; density contour lines in dark red, Spearman's correlation coefficient $\rho$ and associated test's p-value on top left corner. (D) \gls{gsea} \gls{nes} of Hallmark gene sets, ranked by t-statistic, in linear modelled \gls{pl}-isolated \glspl{sg}\cite{somasekharan_g3bp1-linked_2020}. The dashed line corresponds to the least-squares linear regression fit between the two variables (B, C).}
\end{figure}
\noindent
Given the divergent results across \gls{sg} purification strategies, we applied our correction approaches to all remaining datasets that met quality criteria\footnote{Studies were maintained for this analysis if they had appropriate technical controls (i.e., not just SG samples, but also equivalent cellular fractions, empty vectors, etc.) and at least 3 consistent ``\textit{SG}" replicates. Even though most \gls{dc} studies had no appropriate technical controls, we considered that our simulated \gls{sg} samples could serve as such, and thus these studies were maintained.} and compared their transcriptomic profiles. Despite these corrections, substantial discrepancies remained across datasets, independent of purification method (Figure 3.16). The main exception was the set of \gls{dc}-purified \gls{sg} datasets corrected for centrifugation effects, which showed consistency. These observations suggests that transcript-specific sequestration may be minimal or below the detection limits of transcriptome analyses. Alternatively, each purification method may introduce biases that shape the apparent \gls{sg} transcriptome.
\\
\\
\gls{rna} length further complicates interpretation. In \gls{dc}-purified \gls{sg} datasets, correction for centrifugation largely removed the previously reported length bias. In contrast, \gls{pl} datasets showed conflicting patterns: the Somasekharan \gls{etal} 2020 study\cite{somasekharan_g3bp1-linked_2020} shows enrichment of short transcripts, whereas the Ren \gls{etal} 2023 study\cite{ren_profiling_2023} shows enrichment of longer transcripts despite targeting the same \gls{sg} marker protein. Other methods showed additional apparent length biases. \gls{photon}\cite{rajachandran_subcellular_2025} and \gls{fanci}\cite{zhou_rna_2024} isolation approaches were associated with preferential capture of longer transcripts, with \gls{fanci} showing the strongest effect. However, it is important to note that apparent length preferences need not reflect properties specific to any particular \gls{sg} isolation strategy. As mentioned in 3.3.2, \gls{rnaseq} data broadly and recurrently exhibits a sample-specific association between gene length and fold-change estimates, one that persists across commonly used normalisation methods and is likely a general technical artefact of the sequencing technology itself\cite{mandelboum_recurrent_2019}. The length patterns observed across \gls{sg} isolation methods may therefore reflect this general \gls{rnaseq} bias rather than any true enrichment of transcripts of a particular size, although this is merely stated here as an alternative explanation. These considerations highlight the persistent technical challenges involved in reliably characterising \gls{sg} transcriptomes.
\\
\begin{figure}[H]
	\centering
	\includegraphics[width=1\linewidth]{"images/Paper Images/TESE_Fig4.16"}
	\caption[Meta-analysis of SG transcriptomes]{\textbf{Meta-analysis of SG transcriptomes.} Heatmap of Spearman's correlations between transcriptomic profiles of corrected/controlled \glspl{sg} from different studies and \gls{rna} length (\gls{nt}). \gls{sg} purification method of each study is given below the heatmap.}
	\label{fig:Figure4.16}
\end{figure}
\noindent
Although individual gene enrichment varied substantially across studies, we used \gls{gsea} to identify broader regulatory trends. No pathways were consistently enriched or depleted across all SG datasets (Figure 3.17A).
\\
\\
Nevertheless, some consistent signals emerged. During our individual analysis of each dataset, some exhibited enrichment of mitochondrially encoded genes in the putative SG. To investigate this further, we performed GSEA across all datasets using a gene set containing the thirty-six mitochondrially encoded transcripts obtained from Ensembl\cite{cunningham_ensembl_2022}. As a control, we analysed nuclear-encoded transcripts localised to mitochondria, selecting top-scoring genes from the RNALocate database\cite{wu_rnalocate_2025} (Figure 3.17B). Among SG datasets, two studies showed strong depletion of mitochondrially encoded transcripts, whereas four displayed strong enrichment ($|NES| > 2$). No consistent enrichment was observed for nuclear-encoded mitochondrial genes ($|NES| > 2$). To assess the specificity of these signals, we repeated the analysis using 100 random gene sets matched in size to each mitochondrial group, thereby deriving an empirical \gls{fpr}. This yielded \gls{fpr} $<$0.05 in eight mito-encoded datasets, five enriched and three depleted, as well as in one nuclear-encoded depleted dataset. The opposing directions of enrichment and depletion across datasets may reflect genuine variation in experimental conditions, such as cell line, stress type, and duration, rather than noise, as the specificity of the signal relative to size-matched random gene sets argues against a chance finding. Notably, the absence of a consistent signal for nuclear-encoded mitochondrial genes suggests that this effect is specific to mitochondrially encoded transcripts rather than a general property of mitochondria-associated biology. Collectively, these results suggest that mitochondrially encoded transcripts may be dynamically regulated in the context of SG formation in a condition-dependent manner, warranting further investigation.
\\
\\
Intriguingly, this pattern was not easily reconciled with known transcript properties. Mitochondrial transcripts are generally short, yet datasets from Ren \gls{etal}. 2023\cite{ren_profiling_2023} and Rajachandran \gls{etal} 2025\cite{rajachandran_subcellular_2025}, which already showed a strong bias toward longer RNAs, also displayed enrichment of these short mitochondrial transcripts. Moreover, the Kudrin \gls{etal} 2024 study\cite{kudrin_n4-acetylcytidine_2024}, which employed a centrifugation-corrected method, also showed mitochondrial enrichment. Although several other DC datasets showed similar enrichment prior to correction, this signal disappeared after centrifugation correction. Nevertheless, we opted to integrate all of these datasets to uncover whether these or other genes were preferentially enriched.
\\
\\
To derive a consensus SG enrichment signal, we averaged t-statistics across all these SG-transcriptome datasets (Figure 3.17C). Although mito-encoded RNAs are typically enriched, they are not among the most strongly enriched transcripts overall. Among the top-ranked transcripts are several of uncertain biological relevance in this context\cite{cunningham_ensembl_2022}, including non-coding or pseudogenic transcripts such as \textit{MIR4435-2HG} (a \gls{mirna} host gene), \textit{TMEM147-AS1} (a \gls{lncrna}), \textit{CRYBB2P1} (a pseudogene), and \textit{FAM86DP} (another pseudogene).
\\
\\
GSEA of this consensus SG signal (with additional control with the aforementioned \gls{fpr} method) revealed enrichment of \gls{tnfa} signalling via \gls{nfkb} (Figure 3.17D), potentially linked to the attenuation of inflammatory signalling upon SG assembly\cite{paget_stress_2023}, together with depletion of several metabolic pathways, including glycolysis, MYC targets, and oxidative phosphorylation, highlighting biological processes potentially associated with SG-mediated transcript sequestration.

\begin{figure}[H]
	\centering
	\includegraphics[width=1\linewidth]{"images/Paper Images/TESE_Fig4.17"}
	\caption[Consensus RNA Enrichment Across SG Transcriptomes]{}
	\label{fig:Figure4.17}
\end{figure}

\clearpage

\begin{figure}[t]
	\captionsetup{labelformat=adja-page}
	\ContinuedFloat
	\caption[]{\textbf{Mitochondrial RNA Enrichment Across SG Transcriptomes.} (A) (Left) NES from GSEA using Hallmark gene sets, based on genes ranked by t-statistic from differential gene expression between corrected/controlled SGs and WCs across multiple studies. (Right) Paired half-violin plots showing the distribution of NES values obtained from GSEA using 1000 size-matched random gene sets (violins), for mitochondrially encoded (darker shades) or nuclear-encoded mitochondrially localised genes (lighter shades), across the transcriptomic profiles of corrected/controlled SGs from different studies. The observed NES for each gene set is shown as a large dot; filled dots indicate empirical FPR $<$ 0.05. Numbers adjacent to each dot indicate the percentile of the observed NES within the corresponding null distribution. Colour of each violin indicates SG purification method, as indicated above the image. (B) Average t-statistic enrichment plotted against variance across SG-isolated datasets; transcripts as circles, coloured green if mitochondrial encoded, blue if ``strongly depleted" (Variance $<$ 150 and Average Enrichment $<$ -9), and red if ``strongly enriched" (Variance $<$ 150 and Average Enrichment $>$ 7.5) (C) GSEA NES of Hallmark gene sets ranked by the average t-statistic from the same datasets; only gene sets with adjusted p-value $<$0.05 and empirical FPR $<$ 0.05 based on 100 random gene sets of matched size are shown.}
\end{figure}
\noindent